% concatenated from mainTensors.tex
\documentclass[a4paper,12pt,reqno]{amsart}
\usepackage{latexsym}
\usepackage{amssymb} % \mathbb
\usepackage{mathrsfs}
\usepackage{amsmath}
%\usepackage{mathabx}
\usepackage{latexsym}

\usepackage{delarray}
\usepackage{amssymb,amsmath,amsfonts,amsthm,mathrsfs}
\usepackage{blkarray}
%\usepackage{MnSymbol}
%\usepackage[notref,notcite]{showkeys} 

 %%%%%%%%%%%%%%%%%%%%%%%%%%
\setlength{\textwidth}{15.2cm}
\setlength{\textheight}{22.7cm}
\setlength{\topmargin}{0mm}
\setlength{\oddsidemargin}{3mm}
\setlength{\evensidemargin}{3mm}
%\setlength{\headsep}{2.9cm}
\setlength{\footskip}{1cm}


\usepackage{hyperref}
%\usepackage{showkeys}
%\usepackage{showlabels}
%\usepackage[backref]{hyperref}
\renewcommand\eqref[1]{(\ref{#1})} %Need with hyperref
%\usepackage{refcheck} %Must come after hyperref
%%%%%%%%%%%%%%%%%%%%%%%%%%%
%\usepackage{amsmath}
%\usepackage{amsfonts} 
%\usepackage{amssymb}
%\usepackage{latexsym}
\hyphenation{ope-rators}
 \newtheorem{thm}{Theorem}[section]
 \newtheorem{cor}[thm]{Corollary}
 \newtheorem{lem}[thm]{Lemma}
 \newtheorem{prop}[thm]{Proposition}
 \theoremstyle{definition}
 \newtheorem{defn}[thm]{Definition}
 \theoremstyle{remark}
 \newtheorem{rem}[thm]{Remark}
 \newtheorem{ex}[thm]{Example}
 \numberwithin{equation}{section}
\newcommand{\Sbm}{\widetilde{\Sigma_M}}
\newcommand{\half}{\frac{1}{2}}
\newcommand{\twothird}{\frac{2}{3}}
\newcommand{\frp}{\frac{1}{p}}
\newcommand{\fraq}{\frac{1}{q}}
\newcommand{\sdual}{\frac{1}{p}+\frac{1}{q}=1}
\newcommand{\ene}{\mathbb{N}}
\newcommand{\bba}{\mathcal{B}}
\newcommand{\LL}{\mathcal{L}}
\newcommand{\noi}{\noindent}
\newcommand{\frn}{\frac{n}{2}}
\newcommand{\fro}{\frac{1}{p_0}}
\newcommand{\frqo}{\frac{1}{q_0}}
\newcommand{\er}{\mathbb{R}}
\newcommand{\ar}{\mathbb{R}}
\newcommand{\ce}{\mathbb{C}}
\newcommand{\zet}{\mathbb{Z}}
\newcommand{\zn}{\mathbb{Z}^n}
\newcommand{\tn}{\mathbb{T}^n}
\newcommand{\T}{\mathbb{T}^1}
\newcommand{\To}{\mathbb{T}}
\newcommand{\ese}{\mathbb{S}^1}
\newcommand{\ern}{{\mathbb{R}}^n}
\newcommand{\ernn}{{\mathbb{R}}^{4n}}
\newcommand{\fou}{\mathcal{F}}
\newcommand{\efee}{\mathcal{F}}
\newcommand{\efem}{\mathcal{F}}
\newcommand{\efeM}{{\mathcal{F}_M}}
\newcommand{\efeG}{{\mathcal{F}_G}}
\newcommand{\efemo}{\mathcal{F}_M\otimes\overline{\mathcal{F}_M}}
%\newcommand{\ernx}{{\mathbb{R}}_{x}^n}}
%\newcommand{\erns}{{\mathbb{R}}_{\xi}^n}}
\newcommand{\erdn}{{\mathbb{R}}^{2n}}
\newcommand{\bi}{\begin{itemize}}
\newcommand{\cinfm}{{C}^{\infty}(M)}
\newcommand{\cinf}{\mathcal{C}^{\infty}}
\newcommand{\ei}{\end{itemize}}
\newcommand{\be}{\begin{enumerate}}
\newcommand{\ee}{\end{enumerate}}
\newcommand{\beq}{\begin{equation}}
\newcommand{\eq}{\end{equation}}
\newcommand{\OO}{\Omega\times\Omega} 
\newcommand{\efet}{{\mathcal{F}}_{\mathbb{T}^n}}
\newcommand{\efez}{{\mathcal{F}}_{\mathbb{Z}^n}}
\newcommand{\lap}{\mathcal{L}_G}
\newcommand{\cdxi}{\ce^{d_{\xi}\times d_{\xi}}}
\newcommand{\subl}{\mathcal{L}}
\newcommand{\eigen}{\lambda^2_{[\xi]}}
\newcommand{\phs}{M\times \ene\times\ene}
%%%%%%% additional definitions %%%%%
\newcommand{\xy}[2]{\ensuremath{\frac{x_{#1}}{y_{#2}}}}

\newcommand{\Dcal}{\mathcal{D}}
\newcommand{\Dp}{\mathcal{D}_p}
\def\p#1{{\left({#1}\right)}}
\def\b#1{{\left\{{#1}\right\}}}
\def\br#1{{\left[{#1}\right]}}
\def\jp#1{{\left\langle{#1}\right\rangle}}
\def\n#1{{\left\|{#1}\right\|}}
\def\abs#1{{\left|{#1}\right|}}
\def\wt#1{{\widetilde{#1}}}

\def\Rep{{{\rm Rep}}}
\def\Op{{{\rm Op}}}
\def\G{{G}}

\def\En{{E_{o}}}

\def\FT{{\mathscr F}}
%\def\Hcal{{\mathcal H}}
\DeclareMathOperator{\Tr}{Tr}
\DeclareMathOperator{\Det}{Det}
\def\diff{\mathrm{diff}}
\def\Diff{\mathrm{Diff}}
\DeclareMathOperator{\rank}{rank}
\DeclareMathOperator{\singsupp}{sing\,supp}
\def\d{\mathrm d}

\def\Gh{{\widehat{G}}}
\def\sumxi{\sum_{[\xi]\in\Gh}}
\def\dxi{{d_\xi}}
\def\sumij{\sum_{i,j=1}^{\dxi}}
\def\HS{{\mathtt{HS}}}
\def\Rn{{{\mathbb R}^n}}
\def\Tn{{{\mathbb T}^n}}
\def\Zn{{{\mathbb Z}^n}}
\def\T{{{\mathbb T}^1}}
\def\N{{{\mathbb N}}}
\def\C{{{\mathbb C}}}
\def\SU2{{{\rm SU(2)}}}
\def\SO3{{{\rm SO(3)}}}
\def\lapsu2{{{\mathcal L}_\SU2}}
\def\lapsu2{{{\mathcal L}_\SU2}}
\def\LapG{{\mathcal L}_G}
\def\whf{\widehat{f}}
\def\whu{\widehat{u}}

\def\rhodiff{\mbox{$\triangle\!\!\!\!\ast\,$}}


\def\Re{{{\rm Re}\,}}
\def\Op{\text{\rm Op}}
\DeclareMathOperator{\Lip}{Lip}
\DeclareMathOperator{\Ob}{O}
%\DeclareMathOperator{\Tr}{Tr}
\DeclareMathOperator{\Hi}{{\bf H}}
\DeclareMathOperator{\Hcal}{\bf H}
\DeclareMathOperator{\Ki}{{\bf K}}
\def\comment#1{\smallskip{\centering\fbox{\parbox{0.9\textwidth}{\sf #1}}\\}\smallskip}


\begin{document} 

%-------------------------------------------------------------------------
% editorial commands: to be inserted by the editorial office
%
%\firstpage{1} \volume{228} \Copyrightyear{2004} \DOI{003-0001}
%
%
%\seriesextra{Just an add-on}
%\seriesextraline{This is the Concrete Title of this Book\br H.E. R and S.T.C. W, Eds.}
%
% for journals:
%
%\firstpage{1}
%\issuenumber{1}
%\Volumeandyear{1 (2004)}
%\Copyrightyear{2004}
%\DOI{003-xxxx-y}
%\Signet
%\commby{inhouse}
%\submitted{March 14, 2003}
%\received{March 16, 2000}
%\revised{June 1, 2000}
%\accepted{July 22, 2000}
%
%
%
%---------------------------------------------------------------------------
%Insert here the title, affiliations and abstract:
%
\title[Schatten-von Neumann classes of tensors ]
 {Schatten-von Neumann classes of  tensors of invariant operators }


%----------Author 1
\author{Julio Delgado}

\address{Universidad del Valle\\
Departamento de Matematicas\\
Calle 13 100-00\\
Cali\\
Colombia}

\email{delgado.julio@correounivalle.edu.co}

%----------Author 2
\author{Liliana Posada}

\address{Universidad del Valle\\
Departamento de Matematicas\\
Calle 13 100-00\\
Cali\\
Colombia}

\email{liliana.posada@correounivalle.edu.co}



%----------Author 2
\author[Michael Ruzhansky]{Michael Ruzhansky}

\address{%
		Ghent University\\
	Department of Mathematics: Analysis, Logic and Discrete Mathematics\\
	Krijgslaan 281, Building S8 \\
	B 9000 Ghent\\
	Belgium}

\email{Michael.Ruzhansky@ugent.be}

\address{% 
	Queen Mary University of London
	School of Mathematical Sciences\\
	Mile End Road\\
	London E1 4NS\\
	United Kingdom}
\email{m.ruzhansky@qmul.ac.uk}

\thanks{The authors are supported by the FWO Odysseus 1 grant G.0H94.18N: Analysis and Partial Differential Equations and by the Methusalem programme of the Ghent University Special Research Fund (BOF) (Grant number 01M01021). Michael Ruzhansky is also supported by FWO Senior Research Grant G011522N and EPSRC grants EP/R003025/2 and EP/V005529. The first and second authors were also supported by Vic Inv Universidad del Valle CI-71352.}

%----------classification, keywords, date
\subjclass[2020]{Primary 47B10; Secondary 47A80}

\keywords{Tensor products, Schatten-von Neumann norms,  eigenvalues}

\date{\today}
%----------additions
%\dedicatory{To my boss}
%%% ----------------------------------------------------------------------
\begin{abstract}
In this work we study Schatten-von Neumann classes of tensor products of invariant operators on  Hilbert spaces. In the first part we  first deduce some spectral properties for tensors of anharmonic oscillators thanks to the knowledge on corresponding Schatten-von Neumann properties. In the second part we specialised on tensors of invariant operators. In the special case where a suitable Fourier analysis  associated to a fixed partition of a Hilbert space into finite dimensional subspaces  is available we also give the corresponding formulae in terms of symbols. We also give a sufficient condition for Dixmier traceability for a class of finite tensors of pseudo-differential operators on the flat torus. 
\end{abstract}

%%% ----------------------------------------------------------------------
\maketitle
%%% ----------------------------------------------------------------------
%\tableofcontents

\section{Introduction}
The notion of finite tensors of Hilbert spaces appears for the first time in the works of  Murray and von Neumann in \cite{Murray}. The more general case of infinite tensors was extended by von Neumann in \cite{Ne}.  A recent account on these tensors can be found in  \cite{weav:book}.\\

The trace class is a fundamental concept in quantum mechanics. A  {\em  density operator} or a {\em statistical operator} is a  positive self-adjoint trace class operator for a corresponding ensemble of quantum systems, thus a special element of $S_1(\Hi)$. On the other hand, when systems interact the space of states takes the form of a tensor of Hilbert space, that is, a tensor of spaces of states.  The Schatten-von Neumann class  $S_2(\Hi)$  is also known as the class of Hilbert-Schmidt operators. In the physics  literature, $S_2(\Hi)$ it is known as the {\em Liouville space} over $\Hi$. When dealing with coupled quantum systems or open quantum systems, the tensor products of Hilbert spaces arise. More specifically, if $\Hi_1, \Hi_2 $ are  Hilbert spaces and the state spaces of respective quantum systems, then the Hilbert tensor product  $\Hi_1\otimes\Hi_2$ is the state space of the corresponding coupled system. Hilbert tensor products of more factors are required  when the number of systems interacting increases.   An important application of Hilbert tensor products and Schatten-von Neumann classes naturally arises in the study of quantum channels in Quantum Information Theory. The use of Hilbert tensors in Quantum Mechanics goes back to the beginnings of its  mathematical foundations.  \\


The Schatten-von Neumann norms have been recently applied for measurements of quantum statistical speed. In \cite{mgas:sp}, one introduced the notion 
 of Schatten-von Neumann speed.  Therein, each Schatten-von Neumann norm is used to define a corresponding quantum  statistical distance which in turn induces a quantum statistical speed. The measurement of a quantum statistical speed quantifies the sensitivity of an initial state with respect to changes of the parameter of a dynamical evolution. The statistical speed of a quantum state can be interpreted as an observable witness for entanglement. Further applications of Schatten-von Neumann speeds can be found in \cite{schit:cc1}.\\

 
%The notion of Hilbert tensors has been studied since the seminal works of  Murray and von Neumann \cite{}, \cite{} \cite{} and Schatten \cite{} as continuation of von Neumann \cite{}. Further developments were carried out by I. Sigal \cite{} \cite{}. Recent research on operators on Hilbert tensor products and their applications  can be found in  \cite{}, \cite{}, \cite{}.\\
 
 
 
 
 
We will apply the Fourier analysis associated to a fixed partition of a Hilbert space into finite dimensional subspaces as introduced in \cite{fjdmr:foum}. From it one can define a notion of global symbol corresponding to a Fourier multiplier. We then consider tensors of Fourier multipliers and derive sharp Schatten-von Neumann properties and trace formulae  on Hilbert tensor products in terms of global symbols. We will also study some consequences for pseudo-differential operators on the torus that can be decomposed as a composition of a multiplication operator and a invariant operator(Fourier multiplier).\\
   

The authors have investigated different Schatten-von Neumann properties in previous works. Sharp sufficient conditions for kernels of integral operators on compact manifolds \cite{dr:suffkernel} and in other settings \cite {dr:intsc}, anharmonic oscillators \cite{anh:cdr},  Schatten-von Neumann classes and trace formulae for Fourier multipliers  on compact groups  \cite{dr13:schatten}, Grothendieck-Lidskii formulas and nuclearity in \cite{dr13a:nuclp}, Schatten-von Neumann properties for Fourier multipliers associated to partitions on Hilbert spaces \cite{fjdmr:foum}. The Dixmier trace and Wodzicki residue have been investigated for pseudo-differential oeprators on compact manifolds and compact Lie groups in \cite{cdc20}, \cite{ckc20}, \cite{cc20}.\\

In Section \ref{sec2} we review some basic definitions on Schatten-von Neumann Ideals and finite tensor products of Hilbert spaces. Section \ref{sec3} is devoted to the case of Schatten-von Neumann properties on finite  tensor products and the special case of the trace class. We deduce spectral properties for tensors of anharmonic oscillators and in particular for the rate of growth of the energy levels of the tensor.  In Section  \ref{sec4} we review the notion of global symbol for  a class of invariant operators with respect to a partition of a Hilbert space into finite dimensional subspaces in the sense of \cite{fjdmr:foum}. Therein we use the notion of the Fourier analysis associated to that kind  of partitions. In Section \ref{sec5}, we apply the notion of global symbol to the setting of  tensors of invariant operators.  We obtain some Schatten-von Neumann properties for tensors of invariant operators on compact Lie groups. At the end of the section we obtain a sufficient condition for Dixmier traceability for a class of finite tensors of pseudo-differential operators on the flat torus with symbols of the form $\sigma(x,j)=a(x)\beta(j)$ defined on $\Tn\times \Zn$. This class of symbols enjoy some special features for the analysis and have been the object of recent interest in the study under the form of the so-called  pseudo-differential neural operators  (cf. \cite{pnetw:ko}).  

%density matrix or density operator***? what about the terminology?*****

 
%In this paper we will establish Plemelj-Smithies formulas for the determinant $\det(I+A)$ when $A$ is invariant and in particular for negative powers of a suitable partial differential operator on the manifold $M$.  ****  
 
\section{Tensor product of Hilbert spaces and Schatten-von Neumann ideals}\label{sec2}
In this section we recall the basic elements of  tensor products of Hilbert spaces and Schatten-von Neumann ideals. \\

Let $\Hi$ be a complex Hilbert space, we denote its inner product by $\langle \cdot,\cdot\rangle_{\Hi}$ or simply by $\langle \cdot,\cdot\rangle$ when there is not room for confusion. Since the main applications related to this work arise from  quantum mechanics, we adopt here the convention on the antilinearity on the first factor (or conjugate homogeneity) of the inner product, i.e., 
\[ \langle a\phi,\psi\rangle=\overline{a}\langle \phi,\psi\rangle,\]
 for all $\phi,\psi\in\Hi, a\in\ce.$\\

We will sometimes also use the Dirac's {\em Bra} and {\em Ket} notation. If $\psi$ is an element of the Hilbert space  $\Hi$ we will say that it is a {\em ket} and write $|\psi\rangle$. If  $\phi$ is an element of the Hilbert space  $\Hi$ we will associate to it an element of $\Hi^*$ that we call the Bra-$\phi$ denoted by  $\langle \phi|$. The  Bra-$\phi$ defines a
 linear form on $\Hi$ given by 
\[\langle \phi|(|\psi\rangle):=\langle \phi,\psi\rangle.\]
The notation can be simplified and justified by the Riesz representation Theorem writing
\[\langle \phi|\psi\rangle=\langle \phi,\psi\rangle.\]

 
The trace class is a crucial object to define several important concepts in the mathematical framework of quantum mechanics, and in particular it is instrumental  for the main applications that  will be considered herein. We now recall its definition. \\


Let $T:\Hi\rightarrow \Hi$ be an operator in $S_1(\Hi)$ and let  $\{\phi_k\}_k$ be any orthonormal basis for the Hilbert space $\Hi$. Then, the series $\sum\limits_{k=1}^{\infty}\langle \phi_k,T\phi_k\rangle_{\Hi}$ is absolutely convergent and the sum is independent of the choice of the orthonormal basis $\{\phi_k\}_k$. Thus, we can define the trace $\Tr (T)$ of any linear operator
$T:\Hi\rightarrow \Hi$ in $S_1(\Hi)$ by $$\Tr (T)=\sum\limits_{k=1}^{\infty}\langle \phi_k,T\phi_k\rangle_{\Hi},$$
where $\{\phi_k:k=1,2,\dots\}$ is any orthonormal basis for $\Hi$. \\

%We note that in particular the trace above defined can be written in %Dirac's notation as 
%\[\Tr (T)=\sum\limits_{k=1}^{\infty}\langle \phi_k|T|\phi_k\rangle_{\Hi},\]
%no


\subsection{Tensor products of Hilbert spaces and operators}
We shall now recall the definition and basic properties of tensor products of Hilbert spaces and operators which are key concepts in Quantum Mechanics. They arise when several quantum systems are involved and interact with each other. 

\begin{defn}
	Let $\Hi_1, \Hi_2$ be  two complex  Hilbert spaces. Let $\phi\in \Hi_1$ and $\psi\in \Hi_2$. We define 
	 $\phi\otimes\psi$ to be the bi-antilinear form on $ \Hi_1\times \Hi_2$ given by 
	 \[\phi\otimes\psi(\phi',\psi')=\langle \phi',\phi \rangle\langle \psi',\psi \rangle ,\]%


%We recall that if $\mathcal{E}_1,\mathcal{E}_2$ are complex inner spaces, there is a unique product on the algebraic tensor product $\mathcal{E}_1 \widetilde{\otimes} \mathcal{E}_2$ such that for all $\phi,\phi' \in \mathcal{E}_1$ and $\psi,\psi' \in \mathcal{E}_2$, we have 
%\[\langle \phi\otimes \psi, \phi'\otimes \psi'\rangle=\langle \phi',\phi %\rangle\langle \psi',\psi \rangle. \]

\vspace{0.5cm}

 
	Let $\mathcal{E}$ be the linear span of terms of the form $\phi\otimes\psi$, and we will denote it by $ \Hi_1\widetilde{\otimes} \Hi_2$. The vector space $\mathcal{E}=\Hi_1\widetilde{\otimes} \Hi_2$ can be endowed with the inner product  defined by 
	\begin{align*}\langle \phi\otimes\psi, \phi'\otimes\psi'\rangle&=\langle \phi,\phi' \rangle\langle \psi,\psi' \rangle\\
	&= \phi'\otimes\psi'(\phi,\psi).
		\end{align*}
  
In fact, if $u=\sum\limits_{i=1}^n \phi_i \otimes \psi_i \in \mathcal{E}$, then the norm $u$ induced for the product $\langle \cdot, \cdot\rangle$ is
\[\left\|\sum\limits_{i=1}^n \phi_i \otimes \psi_i\right\|=\left\{\sum\limits_{i=1}^n\sum\limits_{k=1}^n \langle \phi_i,\phi_k \rangle\langle \psi_i,\psi_k \rangle\right\}^{1/2}.\]
	%and its natural extension to linear combinations.
 
\end{defn}
One can show that the product $\langle \cdot, \cdot\rangle$ on $\mathcal{E}$ is well defined and positive definite.
 The space $(\mathcal{E}, \langle \cdot, \cdot\rangle)$ is then a pre-Hilbert space. 
  It is clear that we can extend the above definition of tensor product  to the case of a finite number of Hilbert spaces $\Hi_1, \Hi_2,\dots, \Hi_n$; getting again a  corresponding pre-Hilbert space  $(\mathcal{E}, \langle \cdot, \cdot\rangle)$.    
  The notion of Hilbert tensors was introduced by Murray and von Neumann in \cite{Murray} for finite tensors. 

 
 %We point out that, at that moment the authors refer to these ones as direct products.
\begin{defn} The completion of $(\mathcal{E}, \langle \cdot, \cdot\rangle)$ is denoted by  $ \Hi_1\otimes\cdots \otimes \Hi_n$. 
	  The space $\Hi_1\otimes\cdots \otimes \Hi_n$ is called the {\em tensor product} of the  complex  Hilbert spaces $\Hi_1,\dots , \Hi_n$.
\end{defn}

We will denote by $\ene$ the set of natural numbers $\{1,2,\dots\}$ and by $\ene_0$ the set of natural numbers including $0$, i.e. $\ene_0 =\{0,1,2,\dots \}.$ The following theorem is well-known and is a basic property for the construction of a basis in tensor products.  
\begin{thm}\label{prhdes} Let $(e_k)_{k\in \mathcal{N}}$, $(f_{\ell})_{\ell\in \mathcal{M}}$ be orthonormal bases of the complex Hilbert spaces $\Hi_1$ and $\Hi_2$ respectively. Then, the family $(e_j\otimes f_{\ell})_{(k,\ell)\in\mathcal{N}\times \mathcal{M}}$ is an orthonormal basis of $\Hi_1\otimes\Hi_2$.
\end{thm}
It is clear that the previous theorem can be extended to finitely many complex Hilbert spaces  $\Hi_1,\dots , \Hi_n$.
\begin{cor}\label{prhdes33} Let $\Hi_1,\dots , \Hi_n$ be complex Hilbert spaces. Let $(e_k^i)_{k\in \mathcal{N}_i}$ be an orthonormal basis of each $\Hi_i$,  $i=1,2,\dots , n$. Then, $(e_{k_1}^1\otimes\cdots \otimes  e_{k_n}^n)_{(k_1,\dots ,k _n)\in\mathcal{N}_1\times\cdots \mathcal{N}_n}$ is an orthonormal basis of  $ \Hi_1\otimes\cdots \otimes \Hi_n$.
\end{cor}

We observe that the Hilbert spaces are not assumed separable in the previous statements. Non-separable Hilbert spaces arise in several situations,  in particular the countable tensor product of separable Hilbert spaces is not separable.  	
See for instance  \cite{gg:qq1}, \cite{gg:qq2} for some recent model proposals in the context of Quantum Field Theory within the non-separable setting. Hereafter we will only consider complex separable Hilbert spaces.
\medskip%
 
 
 We refer the reader to \cite{weid:b1} for a detailed treatment on the basics of Schatten-von Neumann classes of operators between different Hilbert spaces. Regarding the Hilbert tensor products
 the reader can refer to \cite{r-s:vol1}, \cite{prug:b1} for more comprehensive treatments on these topics. However we should point out that some preliminaries regarding tensor products of Hilbert spaces and specially partial traces are not easily accessible from the existing  literature. In particular we have included some basics on countable tensor products of Hilbert spaces.  
 In this regard we refer to \cite{gui:hi}, \cite{Pah:book} for more detailed discussions for some of those properties. 
\subsection{Schatten-von Neumann classess and the trace}
We now recall the notion of Schatten-von Neumann ideals. Let  $\Hi_1,\Hi_2$ be  complex Hilbert spaces. We denote by $\mathcal{L}(\Hi_1, \Hi_2)$ the algebra of bounded linear operators from $\Hi_1$ into $\Hi_2$ endowed with the operator norm. If $\Hi_1=\Hi_2$ we simply write $\mathcal{L}(\Hi)$. A linear compact operator $A:\Hi_1\rightarrow \Hi_2$ is said to belong to the Schatten-von Neumann class $S_r(\Hi_1,\Hi_2)$ for $0<r<\infty$ if 
$$\sum\limits_{n=1}^{\infty}(s_n(A))^r<\infty,$$ 
where $s_n(A)$ denote the singular values of $A$, i.e. the eigenvalues of $|A|=\sqrt{A^*A}$,  
with multiplicities counted. If $1\leq r<\infty$ 
the class $S_r(\Hi_1,\Hi_2)$ becomes a Banach space endowed with the norm 
\[\|A\|_{S_r}=\left(\sum\limits_{n=1}^{\infty}(s_n(A))^r\right)^{\frac{1}{r}}.\]
If $0<r<1$ the identity above only defines a quasi-norm with respect to which $S_r(\Hi_1,\Hi_2)$ is complete. In  the case $\Hi=\Hi_1=\Hi_2$ we simply denote $S_r(\Hi_1,\Hi_2)=S_r(\Hi)$. The class $S_2(\Hi_1,\Hi_2)$ and $S_1(\Hi)$ are usually known as the class of Hilbert-Schmidt operators and the trace class, respectively. In the case of $r=\infty$ we can put $\|A\|_{S_{\infty}}$ to be the operator norm of the bounded operator
$A:\Hi_1\to \Hi_2$. In this case  $S_{\infty}(\Hi_1,\Hi_2)$ is the class of compact operators from 
 $\Hi_1$ into $\Hi_2$ endowed with the operator norm.\\



\section{Schatten-von Neumann properties of Tensors }\label{sec3}
 We start by recalling the definition of tensors of operators and deducing some basic properties on tensors of Schatten-von Neumann operators.
 

\subsection{Tensor products of operators}
We can now define the tensor product of operators. 
\begin{defn} Let $A$ be an operator on $ \Hi_1$, with a dense domain  $Dom A$, and $B$ be a operator on $ \Hi_2$ with dense domain $Dom B$. Let $D$ denote the space  $Dom A\overline{\otimes} Dom B$, that is the space of finite linear combinations of $\phi\otimes\psi$ with $\phi\in Dom A$ and $\psi\in Dom B$. Clearly $D$ is dense in  $ \Hi_1\otimes \Hi_2$. One defines the operator $A\otimes B$ on  $D$ by 
\[(A\otimes B)(\phi\otimes\psi)=A\phi\otimes B\psi\]
and its linear extension to all $D$.
	\end{defn}
It is clear that the definition above can be natually extended to a finite family of densely defined operators.
\begin{prop}
	The operator $A\otimes B$ is well defined. If $A$ and $B$ are closable operators then so is $A\otimes B$. 
\end{prop}
\begin{defn}
	If $A$ and $B$ are closable operators, we call {\em tensor product } of  $A$ by $B$ the closure of $A\otimes B$ and it still denoted by $A\otimes B$. 
\end{defn}
The following proposition follows immediately from the definition of the tensor $A\otimes B$.
\begin{prop}\label{proboundt} If $A$ and $B$ are bounded operators  on $ \Hi_1$ and  $ \Hi_2$, respectively, then $A\otimes B$ is a bounded operator on  $ \Hi_1\otimes \Hi_2$ and we have
	\[\|A\otimes B\|=\|A\|\|B\|.\] 
\end{prop}

A similar property holds for Schatten-von Neumann classes. We also give some addittional  properties regarding the trace of tensors in the following theorem. We recall their proofs for the convenience of the reader since they are not easy to find in the literature. 

\begin{thm}\label{scprt} Let $\Hi_1, \Hi_2$ be  two complex separable Hilbert spaces. Let $A$ and $B$ be bounded operators  on $ \Hi_1$ and  $ \Hi_2$, respectively.
	\begin{enumerate}
		\item If $A$ and $B$ are compact  operators, then $A\otimes B$ is a compact operator on  $ \Hi_1\otimes \Hi_2$. 

\item Let $0<p<\infty$. If $A\in S_p(\Hi_1)$ and $B\in S_p(\Hi_2)$. Then  $A\otimes B\in S_p( \Hi_1\otimes \Hi_2)$. Moreover we have
\[\|A\otimes B\|_p=\|A\|_p\|B\|_p,\]
and for $p=1$, the following formula for the trace holds
	\[\Tr (A\otimes B)=\Tr (A)\Tr (B).\]
\end{enumerate}
\end{thm}
\begin{proof} (1) Since $A$ and $B$ are compact operators we can write
\[A=\sum\limits_{i=1}^{\infty}\alpha_i|u_i\rangle\langle v_i|,\,\, B=\sum\limits_{j=1}^{\infty}\beta_j|t_j\rangle\langle w_j|, \]
where $(u_i)_i, (v_i)_i $ and $(t_j)_j, (w_j)_j $ are orthonormal sequences in $ \Hi_1$ and  $ \Hi_2$ respectively, and $(\alpha_i)_i, (\beta_j)_j$ are real non-negative sequences converging to $0$.\\

Then, we can write 
\beq A\otimes B=\sum\limits_{i,j=1}^{\infty}\alpha_i\beta_j|u_i\otimes t_j\rangle\langle v_i\otimes w_j|.\label{trAtB}\eq
Hence we obtain a polar representation of $A\otimes B$:
\beq A\otimes B=\sum\limits_{i,j=1}^{\infty}\alpha_i\beta_j|v_i\otimes w_j\rangle\langle v_i\otimes w_j|.\label{tenab}\eq
The compactness now follows immediately since $ (v_i\otimes w_j)_{ij}$ is an orthonormal sequence in $ \Hi_1\otimes \Hi_2$ and $(\alpha_i\beta_j)_{ij}$ converges to $0$ as $(i,j)$ goes to $\infty$.\\

(2) By \eqref{tenab} with  $A\in S_p(\Hi_1), B\in S_p(\Hi_2)$, we note that
\[\sum\limits_{i,j=1}^{\infty}(\alpha_i\beta_j)^p=\sum\limits_{i=1}^{\infty}\alpha_i^p\sum\limits_{j=1}^{\infty}\beta_j^p=\|A\|_p^p\|B\|_p^p .\]
Therefore  $A\otimes B\in S_p( \Hi_1\otimes \Hi_2)$ and 
\[\|A\otimes B\|_p=\|A\|_p\|B\|_p.\]

For $p=1$, in order to obtain the formula for the trace, we first note that there exist orthonormal basis  $(u_i)_i, (v_i)_i $ and $(t_j)_j, (w_j)_j $ of $ \Hi_1$ and  $ \Hi_2$ respectively, and complex sequences $(\alpha_i)_i, (\beta_j)_j$ converging to $0$ such that
\[A=\sum\limits_{i=1}^{\infty}\alpha_i|u_i\rangle\langle v_i|\,,\,\,\, B=\sum\limits_{j=1}^{\infty}\beta_j|t_j\rangle\langle w_j|. \]
Hence
\beq A\otimes B=\sum\limits_{i,j=1}^{\infty}\alpha_i\beta_j|u_i\otimes t_j\rangle\langle v_i\otimes w_j|.\label{trAtB2}\eq

  
Since  $ (v_i\otimes w_j)_{ij}$ is an orthonormal basis of $ \Hi_1\otimes \Hi_2$,  by \eqref{trAtB2} we obtain
 \[ (A\otimes B)(v_k\otimes w_l)=\alpha_k\beta_l|u_k\otimes t_l\rangle\]
 
Therefore
 \begin{align*}
 \Tr(A\otimes B)&=\sum\limits_{i,j=1}^{\infty}\langle v_i\otimes w_j, (A\otimes B)(v_i\otimes w_j)\rangle\\ 
 &=\sum\limits_{i,j=1}^{\infty}\langle v_i\otimes w_j, \alpha_i\beta_j(u_i\otimes t_j)\rangle\\
 &=\sum\limits_{i,j=1}^{\infty}\alpha_i\beta_j\langle v_i\otimes w_j, u_i\otimes t_j\rangle\\
  &=\sum\limits_{i,j=1}^{\infty}\alpha_i\beta_j\langle v_i, u_i\rangle\langle w_j,  t_j\rangle\\
  &=\sum\limits_{i=1}^{\infty}\alpha_i\langle v_i, u_i\rangle
  \sum\limits_{j=1}^{\infty}\beta_j\langle w_j,  t_j\rangle\\
 &=\Tr(A)\Tr(B).
\end{align*} 
\end{proof}
The above theorem  has immediate consequences for finite tensors of operators. The following corollary follows arguing by induction on Proposition \ref{proboundt} and Theorem \ref{scprt}.

\begin{cor} \label{corsp} Let $\Hi_1,\dots,\Hi_n$  be complex separable Hilbert spaces.\\
(i)  If $A_j\in \mathcal{L}(\Hi_j)$, then $A_1\otimes\cdots \otimes A_n\in \mathcal{L}(\Hi_1\otimes\cdots\otimes \Hi_n)$
 and 
 \[\|\bigotimes \limits_{j=1}^nA_j\|_{\mathcal{L}(\Hi_1\otimes\cdots\otimes \Hi_n)}=\prod\limits_{j=1}^{n}\|A_j\|_{\mathcal{L}(\Hi_j)}. \]

\noindent (ii) Let  $0<p<\infty$. If $A_1,\dots, A_n$ belong to $S_p(\Hi_1), \dots, S_p(\Hi_n)$ respectively, then $A_1\otimes\cdots \otimes A_n\in S_p(\Hi_1\otimes\cdots\otimes \Hi_n)$. Moreover 
\[\|\bigotimes \limits_{j=1}^nA_j\|_{S_p(\Hi_1\otimes\cdots\otimes \Hi_n)}=\prod\limits_{j=1}^{n}\|A_j\|_{S_p(\Hi_j)}.\]

In particular, for $p=1$ the trace class we additionally have 

\[\Tr\left(\bigotimes \limits_{j=1}^nA_j\right)=\prod\limits_{j=1}^{n}\Tr(A_j).\]	
\end{cor}%


As an example we study  a system of anharmonic oscillators by using recent results  obtained in \cite{anh:cdr}. We consider  anharmonic oscillators  on  $\Rn$  of the form
 \beq A_{k,\ell}=(-\Delta)^{\ell}+|x|^{2k} \label{gahw1},\eq
where $k,\ell$ are integers $\geq 1$. In \cite{anh:cdr}, it was established  that  $(A_{k,\ell}+1)^{-\mu}$ belongs to the Schatten-von Neumann class $S_p(L^2(\Rn))$ provided $\mu>\frac{(k+\ell)n}{2k\ell p} $.

\begin{thm} Let $k_j,\ell_j$ be integers $\geq 1$ for $j=1,\dots, N$. We assume
\[\mu_j>\frac{(k_j+\ell_j)n}{2k_j\ell_j p} ,\]
    for  $j=1,\dots, N$. We denote $\widetilde{A}_j:=( A_{k_j,\ell_j}+1)^{-\mu_j}$ for $j=1,\dots, N$; with $A_{k_j,l_j}$ anharmonic oscillators as in \eqref{gahw1}. Then,  $\widetilde{A}_j\in S_p(L^2(\Rn))$ for each  $j=1,\dots, N$ and 
   \beq\label{anhsch1}\bigotimes \limits_{j=1}^N\widetilde{A}_j\in S_p(L^2(\mathbb{R}^{nN})).\eq
Moreover, the eigenvalues $\lambda_m$ of $\bigotimes \limits_{j=1}^N\widetilde{A}_j$ satisfy the following rate of decay
\[
		 \lambda_m=o(m^{-\frac{1}{p}})\,,\quad \text{as} \quad m \rightarrow \infty\, .
		 \]
Consequently, the eigenvalues (Energy levels) $E_m$ of $\bigotimes \limits_{j=1}^N A_j$ have a growth of order  at least 
\[
		m^{\frac{1}{p}}\,,\quad \text{as} \quad m \rightarrow \infty\, .
		 \]
\end{thm}
\begin{proof} The first consequence follows from Corollary 5.4 (b) of \cite{anh:cdr}. The property \eqref{anhsch1} follows from  Corollary \ref{corsp} and again  Corollary 5.4(b) of \cite{anh:cdr} and the fact that $L^2(\mathbb{R}^{nN})$ and $\bigotimes \limits_{j=1}^N L^2(\mathbb{R}^{n})$ are isomorphic.\\

The rate of decay is a consequence of \eqref{anhsch1} and Corollary 5.6 of \cite{anh:cdr}, and the estimate on the rate of growth of the $E_m$ follows
\end{proof}

\section{Global symbols of invariant operators with respect to partitions of Hilbert spaces}\label{sec4}%
In this section we recall some basic elements of the notion of Fourier multipliers or invariant operators on Hilbert spaces and their corresponding global symbols introduced in \cite{fjdmr:foum}. These notions will be crucial for the study of tensors of invariant operators in the last section.\\%


Given a complex separable Hilbert space $\Hi$, we will  consider a fixed partition of $\Hi$ into a direct sum of finite dimensional subspaces,
$$\Hcal=\bigoplus_{j} H_{j}.$$
One can  associate a notion of invariance relative to such  partition as we will see from the Theorem \ref{THM:inv-rem} below established in \cite{fjdmr:foum} as Theorem 2.1. We recall it in detail with the corresponding notations since it is essential for our discussion. We denote by $d_j$ the dimension of  $H_{j}$ and by  $\{e_{j}^{k}\}_{1\leq k\leq d_{j}}$  an orthonormal basis of $ H_{j}$. In particular we can apply this notion of invariance to  the setting when $M$ is a compact manifold without boundary, $\Hcal=L^{2}(M)$ and $\Hcal^{\infty}=C^{\infty}(M)$. 
 In this particular example  we can
 think of $\{e_{j}^{k}\}$ being an orthonormal basis given by eigenfunctions
 of an elliptic operator on $M$, and $d_{j}$ the corresponding 
 multiplicities. This notion of invariance has been recently applied in the setting of control theory of Cauchy problems on Hilbert spaces \cite{CGDR:Hi}.  
 
\begin{thm}\label{THM:inv-rem}
	Let $\Hcal$ be a complex separable Hilbert space and let $\Hcal^{\infty}\subset \Hcal$ be a dense
	linear subspace of $\Hcal$. Let $\{d_{j}\}_{j\in\N_{0}}\subset\N$ and let
	$\{e_{j}^{k}\}_{j\in\N_{0}, 1\leq k\leq d_{j}}$ be an
	orthonormal basis of $\Hcal$ such that
	$e_{j}^{k}\in \Hcal^{\infty}$ for all $j$ and $k$. Let $H_{j}:={\rm span} \{e_{j}^{k}\}_{k=1}^{d_{j}}$,
	and let $P_{j}:\Hcal\to H_{j}$ be the corresponding orthogonal projection.
	For $f\in\Hcal$, we denote $$\widehat{f}(j,k):=(f,e_{j}^{k})_{\Hcal}$$ and let
	$\widehat{f}(j)\in \ce^{d_{j}}$ denote the column of $\widehat{f}(j,k)$, $1\leq k\leq d_{j}.$
	Let $T:\Hcal^{\infty}\to \Hcal$ be a linear operator.
	Then the following
	conditions are equivalent:
	\begin{itemize}
		\item[(A)] For each $j\in\ene_0$, we have $T(H_j)\subset H_j$. 
		\item[(B)] For each $\ell\in\ene_0$ there exists a matrix 
		$\sigma(\ell)\in\ce^{d_{\ell}\times d_{\ell}}$ such that for all $e_j^k$ 
		$$
		\widehat{Te_j^k}(\ell,m)=\sigma(\ell)_{mk}\delta_{j\ell}.
		$$
		\item[(C)]  If in addition, $e_j^k$ are in the domain of $T^*$ for all $j$ and $k$, then 
		for each $\ell\in\ene_0 $ there exists a matrix 
		$\sigma(\ell)\in\ce^{d_{\ell}\times d_{\ell}}$ such that
		\[\widehat{Tf}(\ell)=\sigma(\ell)\widehat{f}(\ell)\]
		for all $f\in\Hcal^{\infty}.$
	\end{itemize}
	
	The matrices $\sigma(\ell)$ in {\rm (B)} and {\rm (C)} coincide.
	
	The equivalent properties {\rm (A)--(C)} follow from the condition 
	\begin{itemize}
		\item[(D)] For each $j\in\ene_0$, we have
		$TP_j=P_jT$ on $\Hcal^{\infty}$.
	\end{itemize}
	If, in addition, $T$ extends to a bounded operator
	$T\in{\mathscr L}(\Hcal)$ then {\rm (D)} is equivalent to {\rm (A)--(C)}.
\end{thm} 

Under the assumptions of Theorem \ref{THM:inv-rem}, we have the direct sum 
decomposition
\begin{equation}\label{EQ:sum}
\Hcal = \bigoplus_{j=0}^{\infty} H_{j},\quad H_{j}={\rm span} \{e_{j}^{k}\}_{k=1}^{d_{j}},
\end{equation}
and we have $d_{j}=\dim H_{j}.$
The two applications that we will consider will be with $\Hcal=L^{2}(M)$ for a
compact manifold $M$ with $H_{j}$ being the eigenspaces of an elliptic 
pseudo-differential operator $E$, or with $\Hcal=L^{2}(G)$ for a compact Lie group
$G$ with $$H_{j}=\textrm{span}\{\xi_{km}\}_{1\leq k,m\leq d_{\xi}}$$ for a
unitary irreducible representation $\xi\in[\xi]\in\widehat{G}$. The difference
is that in the first case we will have that the eigenvalues of $E$ corresponding
to $H_{j}$'s are all distinct, while in the second case the eigenvalues of the Laplacian
on $G$ for which $H_{j}$'s are the eigenspaces, may coincide.

\begin{defn}\label{inva1}
In view of properties (A) and (C), respectively, an operator $T$ satisfying any of
the equivalent properties (A)--(C) in
Theorem \ref{THM:inv-rem}, will be called an {\em invariant operator}, or
a {\em Fourier multiplier relative to the decomposition
	$\{H_{j}\}_{j\in\N_{0}}$} in \eqref{EQ:sum}.
If the collection $\{H_{j}\}_{j\in\N_{0}}$
is fixed once and for all, we can just say that $T$ is {\em invariant}
or a {\em Fourier multiplier}.

The family of matrices $\sigma$ will be
called the {\em matrix symbol of $T$ relative to the partition $\{H_{j}\}$ and to the
	basis $\{e_{j}^{k}\}$.}
It is an element of the space $\Sigma$ defined by
\begin{equation}\label{EQ:Sigma1}
\Sigma=\{\sigma:\N_{0}\ni\ell\mapsto\sigma(\ell)\in \ce^{d_{\ell}\times d_{\ell}}\}.
\end{equation}
\end{defn}
%A criterion for the extendability of $T$ to ${\mathscr L}(\Hcal)$ in terms
%of its symbol will *******be given in Theorem \ref{L2-abstract}.\\

For $f\in\Hcal$, in the notation of Theorem \ref{THM:inv-rem},
by definition we have
\begin{equation}\label{EQ:ser}
f=\sum_{j=0}^{\infty} \sum_{k=1}^{d_{j}} \widehat{f}(j,k) e_{j}^{k}
\end{equation}
with the convergence of the series in $\Hcal$.
Since $\{e^k_j\}_{j\geq 0}^{1\leq k\leq d_j}$ is a complete orthonormal 
system on  $\Hcal$, for all $f\in \Hcal$ we have the Plancherel formula
\beq \label{EQ:Plancherel}
\|f\|^2_{\Hcal}=\sum\limits_{j=0}^{\infty}\sum\limits_{k=1}^{d_j}|( f,e_j^k)|^2
=  \sum\limits_{j=0}^{\infty}\sum\limits_{k=1}^{d_j}|\widehat{f}(j,k)|^{2}
=\|\widehat{f}\|^{2}_{\ell^2(\N_{0},\Sigma)},
\eq
where we interpret $\widehat{f}\in\Sigma$ as an element of the space
\begin{equation}\label{EQ:aux3}
\ell^2(\N_{0,}\Sigma)=
\{h:\ene_0\rightarrow \prod\limits_d\ce^{d}: h(j)\in \ce^{d_j}\, \mbox{ and }\,\sum\limits_{j=0}^{\infty}\sum\limits_{k=1}^{d_j}|h(j,k)|^2<\infty\}, 
\end{equation}
and where we have written $h(j,k)=h(j)_k$. 
In other words, $\ell^2(\N_{0,}\Sigma)$ is the space of all $h\in\Sigma$ such that
$$
\sum\limits_{j=0}^{\infty}\sum\limits_{k=1}^{d_j}|h(j,k)|^2<\infty.
$$
We endow  $\ell^2(\N_{0},\Sigma)$ with the norm
\begin{equation}\label{EQ:aux4}
\|h\|_{\ell^2(\N_{0,}\Sigma)}:=\left(\sum\limits_{j=0}^{\infty}\sum\limits_{k=1}^{d_j}|h(j,k)|^2\right)^{\half}.
\end{equation}

We note that the matrix symbol $\sigma(\ell)$ depends 
not only on the partition \eqref{EQ:sum} but also
on the choice of the orthonormal basis.
Whenever necessary, we will indicate the dependance of $\sigma$ on the orthonormal 
basis by writing $(\sigma,\{e_j^k\}_{j\geq 0}^{1\leq k\leq d_j} )$ and we also will refer to 
$(\sigma,\{e_j^k\}_{j\geq 0}^{1\leq k\leq d_j} )$ as the {\em symbol} of $T$. 
Throughout this  section the orthonormal basis will be fixed and unless there is some 
risk of confusion the symbols will be denoted simply by $\sigma$.  
In the invariant language,  
we have that the transpose of the symbol,
$\sigma(j)^{\top}=T|_{H_{j}}$ is just the restriction of
$T$ to $H_{j}$, which is well defined in view of the property (A).

We will also sometimes 
write $T_{\sigma}$ to indicate that $T_{\sigma}$ is an operator corresponding to the 
symbol $\sigma $. It is clear from the definition that invariant operators are 
uniquely determined by their symbols. Indeed, if $T=0$ we obtain 
$\sigma=0$  for any choice of an orthonormal basis.  
Moreover, we note that by taking $j=\ell$ in (B) of Theorem \ref{THM:inv-rem} we obtain
the formula for the symbol:
\beq\label{symbinv}
\sigma(j)_{mk}=\widehat{Te_j^k}(j,m),
\eq
for all $1\leq k,m\leq d_j$. The formula (\ref{symbinv}) furnishes 
an explicit formula for the symbol in terms of the operator and the orthonormal basis. 
The definition of Fourier coefficients tells us that for invariant operators
we have 
\beq\label{symbinv2}
\sigma(j)_{mk}=({Te_j^k},e_j^m)_{H}.
\eq
In particular,  for the identity operator $T=I$ we have $\sigma_{I}(j)=I_{d_{j}}$,
where $I_{d_{j}}\in \C^{{d_{j}}\times {d_{j}}}$ is the identity matrix.

Let us now indicate a formula relating symbols 
with respect to different orthonormal bases. 
If $\{e_{\alpha}\}$ and $\{f_{\alpha}\}$ are orthonormal bases of 
$\Hcal$, we consider the unitary operator $U$ determined by 
$U(e_{\alpha})=f_{\alpha}$. Then we have
\[
%\begin{multline*} \label{EQ:equiv-sim}
(Te_{\alpha}, e_{\beta})_{\Hcal}=(UTe_{\alpha}, Ue_{\beta})_{\Hcal}
=(UTU^*Ue_{\alpha}, Ue_{\beta})_{\Hcal}
=(UTU^*f_{\alpha}, f_{\beta})_{\Hcal}.
%\end{multline*}
\]
Thus, if $(\sigma_{T}, \{e_{\alpha}\})$ denotes the symbol of $T$ with respect to the 
orthonormal basis $\{e_{\alpha}\}$ and $(\sigma_{UTU^*}, \{f_{\alpha}\})$ 
denotes the symbol of $UTU^*$ with respect to the orthonormal basis $\{f_{\alpha}\}$ 
we have obtained the relation
\beq\label{difsymb} (\sigma_{T}, \{e_{\alpha}\})=({\sigma_{UTU^*}}, \{f_{\alpha}\}).\eq
Thus, the equivalence relation of bases $\{e_{\alpha}\}\sim  \{f_{\alpha}\}$ given by 
a unitary operator $U$ induces
the equivalence relation on the set $\Sigma$ of symbols given by 
\eqref{difsymb}. In view of this,
we can also think of the symbol being independent of a choice of basis, as an element of the space
$\Sigma/\sim$ with the equivalence relation given by
\eqref{difsymb}.

\section{ Symbols on Hilbert Tensor Products}\label{sec5}
We will now apply the above notion of global symbols to the setting of Hilbert tensor products. First we will see the decomposition of Hilbert tensors as direct sums so that we can apply our global symbols. Secondly, we deduce some  Schatten-von Neumann properties for tensors of invariant operators.  \\

The proof of the next lemma follows from Theorem \ref{prhdes}. 
\begin{lem}\label{lesh2x} Let $\Hi_1, \Hi_2$  be two complex separable Hilbert spaces decomposed as direct sums in the form
\[\Hi_1=\bigoplus\limits_{j=1}^{\infty}\Hi_{1,j}\, ,\,\, \Hi_2=\bigoplus\limits_{k=1}^{\infty}\Hi_{2,k}. \]
Then, the Hilbert space $\Hi_1\otimes\Hi_2$ can be written in the form 
\begin{equation}\label{lempart2h}
    \Hi_1\otimes\Hi_2=\bigoplus\limits_{j,\,k=1}^{\infty}(\Hi_{1,j}\otimes\Hi_{2,k})
    \end{equation}
\end{lem}
As a consequence we obtain formulae for the special case of Hilbert tensor of Fourier multipliers or invariant operators relative to fixed decompositions of the Hilbert factors in the sense of the Definition \ref{inva1}. This theorem shows that the notion of invariance herein considered is well behaved with respect to tensors. First we note that the dimension  of $\Hi_{1,j}\otimes\Hi_{2,k}$ is $d_jd_k$, and a typical element of the tensor of the corresponing orthonormal bases is of the form $e_{1,j}^p\otimes e_{2,k}^q$.
\begin{thm}\label{lesh2x2} Let $\Hi_1, \Hi_2$  be two complex separable Hilbert spaces decomposed as direct sums in the form
\[\Hi_1=\bigoplus\limits_{j=1}^{\infty}\Hi_{1,j}\, ,\,\, \Hi_2=\bigoplus\limits_{k=1}^{\infty}\Hi_{2,k}. \]

Let us consider two invariant operators 
\[\Op(\sigma_1):\Hi_1\rightarrow\Hi_1,\,\, \Op(\sigma_2):\Hi_2\rightarrow\Hi_2,\]
corresponding to symbols $\sigma_1, \sigma_2$ in the sense of Definition \ref{inva1}. Then, the corresponding tensor product of operators $\Op(\sigma_1)\otimes\Op(\sigma_2):\Hi_1\otimes\Hi_2\rightarrow \Hi_1\otimes\Hi_2$, is invariant with respect to the partition \eqref{lempart2h} and  we have
\[ \Op(\sigma_1)\otimes\Op(\sigma_2)=\Op(\sigma_1\otimes\sigma_2),\]
where 
\[(\sigma_1\otimes\sigma_2)(j,k)=\sigma_1(j)\otimes\sigma_2(k).\]
\end{thm}
\begin{proof} The fact that $\Op(\sigma_1)\otimes\Op(\sigma_2)$ is invariant follows by verifying the Condition (A) in  Theorem \ref{THM:inv-rem}. Indeed, since 
\[\Hi_1=\bigoplus\limits_{j=1}^{\infty}\Hi_{1,j}\, ,\,\, \Hi_2=\bigoplus\limits_{k=1}^{\infty}\Hi_{2,k} \]
and $\Op(\sigma_1)(\Hi_{1,j})\subset \Hi_{1,j},\,\, $
$ \Op(\sigma_2)(\Hi_{2,j})\subset \Hi_{2,j} $. We obtain %
\begin{align*}
 \left(\Op(\sigma_1)\otimes\Op(\sigma_2) \right)(\Hi_{1,j}\otimes\Hi_{2,k})&\subset \Op(\sigma_1)(\Hi_{1,j})\otimes \Op(\sigma_2)(\Hi_{2,k}) \\
 &\subset \Hi_{1,j}\otimes\Hi_{2,k}.
\end{align*}%
This shows the invariance of $\Op(\sigma_1)\otimes\Op(\sigma_2).$ In order to obtain the symbol of $\Op(\sigma_1)\otimes\Op(\sigma_2)$, we first observe that for an element $e_{1,j}^p\otimes e_{2,k}^q$ of the orthonormal basis of $\Hi_1\otimes \Hi_2$ we have
\begin{equation}\label{prinrh1}
(\Op(\sigma_1)\otimes\Op(\sigma_2))(e_{1,j}^p\otimes e_{2,k}^q)=\Op(\sigma_1)(e_{1,j}^p)\otimes\Op(\sigma_2)(e_{2,k}^q).
\end{equation}
We now use the identity \eqref{symbinv2} and see that  
\begin{align*}
\left\langle \Op(\sigma_1)(e_{1,j}^p)\otimes\Op(\sigma_2)(e_{2,k}^q), 
 e_{1,j}^r\otimes e_{2,k}^s\right\rangle &=\left\langle \Op(\sigma_1)(e_{1,j}^p), e_{1,j}^r\right\rangle\left\langle \Op(\sigma_2)(e_{2,k}^q), e_{2,k}^s\right\rangle\\
 &=\sigma_1(j)_{(r,p)}\sigma_2(k)_{(s,q)}.
 \end{align*}
By writing $T=\Op(\sigma_1)\otimes\Op(\sigma_2)$, we have that the symbol $\sigma_T$ of $T$, is given by
\[\sigma_T(j,k)_{(r,s),(p,q)}=\sigma_1(j)_{(r,p)}\sigma_2(k)_{(s,q)}.\]
Therefore
\[(\sigma_1\otimes\sigma_2)(j,k)=\sigma_1(j)\otimes\sigma_2(k),\]
completing the proof.
\end{proof}
Theorem \ref{lesh2x2} can of course  be extended to the tensor of finitely many invariant operators. 
\begin{cor} \label{cortenx}  Let $\Hi_1, \Hi_2,\dots, \Hi_n$  be    complex separable Hilbert spaces 
 and assume each one can be decomposed as direct sums in the form
\[\Hi_i=\bigoplus\limits_{j=1}^{\infty}\Hi_{i,j}. \]
Let us consider invariant operators 
\[\Op(\sigma_i):\Hi_i\rightarrow\Hi_i, \]
with corresponding  symbols $\sigma_i$ for $i=1,2,\dots, n$ in the sense of Definition \ref{inva1}. Then, the corresponding tensor product of operators 
\[\bigotimes\limits_{i=1}^{n} \Op(\sigma_i):\bigotimes\limits_{i=1}^{n}\Hi_i\rightarrow \bigotimes\limits_{i=1}^{n}\Hi_i,\] is invariant with respect to the partition
\begin{equation}\label{part29k}
\bigotimes\limits_{i=1}^{n}\Hi_i= \bigoplus\limits_{j_1,j_2,\dots,j_n=1}^{\infty}(\Hi_{1,j_1}\otimes\Hi_{2,j_2}\otimes\cdots\otimes \Hi_{n,j_n}).
\end{equation}

The symbol of the tensor satisfies 
\[ \bigotimes\limits_{i=1}^{n}\Op(\sigma_i)=\Op\left(\bigotimes\limits_{i=1}^{n}\sigma_i\right),\]
where
\[\left(\bigotimes\limits_{i=1}^{n}\sigma_i\right)(j_1,j_2,\dots,j_n)=\bigotimes\limits_{i=1}^{n}\sigma_i(j_i).\]
 \end{cor}
\begin{proof} The fact that \eqref{part29k} holds follows by induction on Lemma \ref{lesh2x}: The invariance of $\bigotimes\limits_{i=1}^{n} \Op(\sigma_i)$ with respect to such partition follows by induction applying Theorem \ref{lesh2x2} as well as the corresponding formula for the symbol. 
\end{proof}
We now give some consequences of Corollary \ref{corsp}, the corollary above and the Theorem  2.5 of \cite{fjdmr:foum} in the setting of Schatten-von Neumann classes.
\begin{cor} \label{invSchatten0} Let  $0<p<\infty$ and   $\Hi_j$  complex separable Hilbert spaces for $j=1, 2, \dots \, n$. If $A_j\in S_p(\Hi_j)$ are invariant  for $j=1, 2, \dots \, n$, then  
\[\bigotimes \limits_{j=1}^{n}A_j\in S_p\left(\bigotimes \limits_{j=1}^{n}\Hi_j\right).\]
Moreover
\[\|\bigotimes \limits_{j=1}^{n}A_j\|_{S_p\left(\bigotimes \limits_{j=1}^{n}\Hi_j\right)}=\prod\limits_{j=1}^{n}\|A_j\|_{S_p(\Hi_j)}=\prod\limits_{j=1}^{n}\left(\sum\limits_{\ell=1}^{\infty}\|\sigma_j(\ell)\|_{S_p(\Hi_{\ell})}^p \right)^{\frac{1}{p}}.\]
In the case of the  trace class $(p=1)$, we additionally have 
\begin{equation}
\Tr\left(\bigotimes \limits_{j=1}^n A_j\right)=\prod\limits_{j=1}^{n}\Tr(A_j)=\prod\limits_{j=1}^{n}\left(\sum\limits_{\ell=1}^{\infty}\Tr(\sigma_j(\ell) )\right).\label{trinf2h}\end{equation}	 
\end{cor}
%Then, the Hilbert space $\Hi_1\otimes\Hi_2$ can be written in the form 
%\[\Hi_1\otimes\Hi_2=\bigoplus\limits_{j,\,k=1}^{\infty}(\Hi_{1,j}\otimes\Hi_{2,k})\]

%**** \[\sigma(j,k):\Hi_{1,j}\otimes\Hi_{2,k}\rightarrow \Hi_{1,j}\otimes\Hi_{2,k}\]









%consider two complex separable Hilbert spaces $\Hi_1, \Hi_2$ decomposed as direct sums in the form
%\[\Hi_1=\bigoplus\limits_{j=1}^{\infty}\Hi_{1,j},\,\, \Hi_2=\bigoplus\limits_{k=1}^{\infty}\Hi_{2,k} \]

%******We dont need invariant here

%In particular, for $p=1$ the trace class we additionally have 
%\[\Tr\left(\bigotimes \limits_{j=1}^n A_j\right)=\prod\limits_{j=1}^{n}\int\limits_{G}\sum\limits_{\xi}\sigma_j(x,\xi)dx.\]	
As an example we derive formulae for Schatten-von Neumann norms for tensors of negative powers $(I-\mathcal{L}_{\SU2})^{-\frac{\alpha}{2}}$ of the Laplacian on $\SU2$. By applying  Corollary \ref{invSchatten0} above and Corollary 4.5 of \cite{dr13:schatten}, we have: %

\begin{cor}\label{COR:su2-Lap}%
Let $0<p<\infty$ and  $\alpha,\, \beta> \frac{3}{p}$ . Then, the tensor 
\[(I-\lapsu2)^{-\frac{\alpha}{2}}\otimes (I-\lapsu2)^{-\frac{\beta}{2}}\]
is invariat and belongs to $S_p(L^2(\SU2\times\SU2))$.\\

Moreover, the norm of the tensor $(I-\lapsu2)^{-\frac{\alpha}{2}}\otimes (I-\lapsu2)^{-\frac{\beta}{2}}$ satisfies
\begin{align*}
\|(I-\lapsu2)^{-\frac{\alpha}{2}}\otimes (I-\lapsu2)^{-\frac{\beta}{2}}\|_{S_p}&=\|(I-\lapsu2)^{-\frac{\alpha}{2}}\|_{S_p}\|(I-\lapsu2)^{-\frac{\beta}{2}}\|_{S_p} 
\end{align*}
\[=\left( \sum\limits_{\ell\in\frac{1}{2}\ene_0}(2\ell+1)^2(1+\ell(\ell+1))^{-\frac{\alpha}{2}p} \sum\limits_{l\in\frac{1}{2}\ene_0}(2l+1)^2(1+l(l+1))^{-\frac{\beta}{2}p}\right)^{\frac{1}{p}}.\]

\end{cor}

\begin{proof} By  Corollary 4.5 of \cite{dr13:schatten} we have that the invariant operators 
$(I-\lapsu2)^{-\frac{\alpha}{2}}$ and $(I-\lapsu2)^{-\frac{\beta}{2}}$ belong to $S_p(L^2(\SU2))$. From Corollary \ref{invSchatten0} we have that the tensor product\\ $(I-\lapsu2)^{-\frac{\alpha}{2}}\otimes (I-\lapsu2)^{-\frac{\beta}{2}}$ is invariant and  belongs to $S_p(L^2(\SU2)\otimes 
 L^2(\SU2))$. \\
 
 Since $L^2(\SU2)\otimes L^2(\SU2)$ is isomorphic to $L^2(\SU2\times\SU2)$, 
 we conclude the proof of the first part.\\

On the other hand, the formula for the Schatten-von Neumann norm of each factor follows from the formulae for the global symbol of  $\lapsu2$ (see \cite{dr13:schatten}).     
\end{proof}
Similar conclusions may be drawn for negative powers of the sub-Laplacian and trace formulae.
Other examples can include a family of `Schr{\"o}dinger operators' on $\SO3$:
\[\mathcal H_{\gamma}=iD_3-\gamma(D_1^2+D_2^2),\]
for a parameter $0<\gamma<\infty,$ and where we give fix three invariant vector fields $D_1, D_2, D_3$ on $\SO3$ corresponding to the derivatives with
respect to the Euler angles. We refer to \cite[Chapter 11]{rt:book} for the explicit formulae.\\

In that case the matrix-symbol of $I+\mathcal H_{\gamma}$ is given by 
\begin{equation}\label{EQ:SO3-schor}
\sigma_{I+\mathcal H_{\gamma}}(\ell)_{mn}=(1+m-\gamma m^2+\gamma\ell(\ell+1))\delta_{mn},\quad m,n\in\mathbb Z,\;
-\ell\leq m,n\leq\ell,
\end{equation}
where as before $\delta_{mn}$ is the Kronecker delta, and we let $m,n$ run from
$-\ell$ to $\ell$ rather than from $0$ to $2\ell+1$.\\

In this last part of the paper we deduce some consequences for a class of periodic pseudodifferential operators, or pseudodifferential operators on the flat torus. We will be mainly focused on those ones with symbols of the form $\sigma(x,j)=a(x)\beta(j)$ defined on $\Tn\times \Zn$. This type of operators has been of recent interest for computational purposes with neural networks  (cf. \cite{pnetw:ko}), where families of operators of this kind play an essential role. Some recent results regarding singular traces for pseudo-differential operators on the flat torus can also be found in  \cite{Piet:trace2}.\\

We will consider a family of symbols $\sigma_m$ of the form  $\sigma_m(x,j)=a_m(x)\beta_m(j)$, where $a_m:\T\rightarrow\er$ is a nonnegative measurable function for $m=1,\dots, N$; $\beta_m:\zet\rightarrow \er$ 
is a nonegative function for $m=1,\dots, N$ such that 
\[ (A1)\,\,\,\|a_m\|_{L^{\infty}(\T)}<\infty,\,\, \sum\limits_{j\in\zet}\beta_m(j)^p<\infty\]
  for all $1<p<\infty$ and for all  $m=1,\dots, N$.\\

In the following theorem, $\Tr_{\omega}$ denotes the Dixmier trace (cf.\cite{Connes}), and we give a sufficent condition for the existence of Dixmier trace for finite tensors of the operators above described.

\begin{thm} Let $A_m$ be pseudodifferential operators with symbol of the form $\sigma_m(x,j)=a_m(x)\beta_m(j)$ for for $m=1,\dots, N$. We assume that they satisfy $(A1)$ and that the limit
\beq\label{dix1} \lim\limits_{p\rightarrow 1^{+}}(p-1)\prod\limits_{m=1}^{N}\|a_m\|_{L^{\infty}(\T)}^p\|\beta_m\|_{\ell^{p}(\zet)}^p\eq
exists. Then, 
\beq \label{dix2}
\lim\limits_{p\rightarrow 1^{+}}(p-1)\Tr\left(\left(\bigotimes_{m=1}^N A_m\right)^p \right) 
\eq
exists, and 
\beq \label{dix3}\Tr_{\omega}\left(\bigotimes_{m=1}^N A_m\right)=\lim\limits_{p\rightarrow 1^{+}}(p-1)\Tr\left(\bigotimes_{m=1}^N A_m^p\right) .\eq
\end{thm}
\begin{proof} We first note that since each $A_m$ is a positive definite operator on $L^{2}(\T)$, then so is 
$\bigotimes_{m=1}^N A_m$ on $\bigotimes_{m=1}^N L^{2}(\T)$, and 
\[\left(\bigotimes_{m=1}^N A_m\right)^p=\bigotimes_{m=1}^N A_m^p,\]
for all $1<p<\infty$. \\

We now note that, $A_m^p$ is a trace class operator for all $1<p<\infty$. Indeed, since $A_m$ is a positive definite operator belonging to the Schatten-von Neumann class $S_p$ as we will see below and 
\[\|A_m^p\|_{S_1}=\|A_m\|_{S_p}^{p}. \]
The operator $A_m$ can be written as a composition of a multiplication operator that is bounded with a Fourier multiplier that is trace class since the symbol of $A_m$ is   $a_m(x)\beta_m(j)$. We observe that
\begin{eqnarray*}
\|A_m\|_{S_p}\leq &\|T_{a_m}\|_{\mathcal{L}(L^2(\T))}\|\beta_m(D)\|_{S_p}\\
 \leq& \|a_m\|_{L^{\infty}(\T)}\|\beta_m(D)\|_{S_p}\\
 =&  \|a_m\|_{L^{\infty}(\T)}\left(\sum\limits_{j\in\zet}\beta_m(j)^p\right)^{\frac{1}{p}}
\end{eqnarray*}
where we have denote by $T_{a_m}$ the multiplication operator associated to $a_m$ and  by $\|\cdot\|_{\mathcal{L}(L^2(\T))}$ the operator norm with respect to the $L^2(\T)$ norm. From this and the assumption in the theorem we get that  $A_m^p$ is a trace class operator for all $1<p<\infty$. 


Then, we have 
\begin{eqnarray*} \Tr\left(\left(\bigotimes_{m=1}^N A_m\right)^p\right)=&\Tr\left(\bigotimes_{m=1}^N A_m^p\right)\\ 
    =& \prod\limits_{m=1}^{N}\Tr(A_m^p)\\
    =& \prod\limits_{m=1}^{N}\|A_m\|_{S_p}^{p}\\
    \leq &  \prod\limits_{m=1}^{N} \|a_m\|_{L^{\infty}(\T)}^p\sum\limits_{j\in\zet}\beta_m(j)^p.
\end{eqnarray*}
From this last inequality and \eqref{dix1}, we can conclude  that the limit  \eqref{dix1} exists and 
\eqref{dix3} follows from  the Connes-Moscovici formula (cf. Proposition 4, \cite{Connes}).
\end{proof}


\noindent {\bf{Declarations}}\\


\noindent $\bullet$ Our manuscript has not associated data.\\


\noindent $\bullet$ No conflict of interest/Competing interests 

	



%\bibliographystyle{alphaabbr}

%\bibliography{bib-Delgado-29-04-14}

\begin{thebibliography}{LYLW20}

%\bibitem[Att10]{Attal:book}
%S.~Attal.
%\newblock {\em {L}ectures in {Q}uantum {T}heory}.
%\newblock Lecture notes, 2010.


\bibitem[CDGR23]{CGDR:Hi}
D. Cardona, J. Delgado, B. Grajales, M. Ruzhansky.
\newblock {C}ontrol of the Cauchy problem on Hilbert spaces: A global approach via symbol criteria. {\em Comm. Pure Appl. Anal.}, 22, 3295-3329, 2023

\bibitem[CDC20]{cdc20} D. Cardona, C. Del Corral. The Dixmier trace and the non-commutative residue for multipliers on compact manifolds. In: Georgiev V., Ozawa T., Ruzhansky M., Wirth J. (eds) Advances in Harmonic Analysis and Partial Differential Equations. Trends in Mathematics. Birkh\"auser, Cham, (2020).

\bibitem[CKC20]{ckc20} Cardona, D., Kumar, V., Del Corral, C. Dixmier traces for discrete pseudo-differential operators, J. Pseudo-Differ. Oper. Appl., Vol. 11, 647--656, 2020.

\bibitem[CC20]{cc20}  Cardona, D., Del Corral, C. The Dixmier trace and the Wodzicki residue for pseudo-differential operators on compact manifolds, Rev. Integr. Temas. Mat., Vol. 38 (1), 67-79, 2020.


\bibitem[CDR21]{anh:cdr}
M.~Chatzakou, J.~Delgado, and M.~Ruzhansky.
\newblock {O}n a class of anharmonic oscillators.
\newblock {\em J. Math. Pures Appl.}, 153:1--29, 2021.

\bibitem[CO94]{Connes} Connes, A. 
\newblock {N}oncommutative Geometry.
\newblock {\em Academic Press.}, 1994.



\bibitem[DR14a]{dr13a:nuclp}
J.~Delgado and M.~Ruzhansky.
\newblock ${L}^p$-nuclearity, traces, and {G}rothendieck-{L}idskii formula on
  compact {L}ie groups.
\newblock {\em J. Math. Pures Appl.}, 102:153--172, 2014.

\bibitem[DR14b]{dr:suffkernel}
J.~Delgado and M.~Ruzhansky.
\newblock Schatten classes on compact manifolds: {K}ernel conditions.
\newblock {\em {J}. {F}unct. {A}nal.}, 267:772--798, 2014.

\bibitem[DR17]{dr13:schatten}
J.~Delgado and M.~Ruzhansky.
\newblock {S}chatten classes and traces on compact groups.
\newblock {\em Math. Res. Letters}, 24:979--1003, 2017.

\bibitem[DR18]{fjdmr:foum}
J.~Delgado and M.~Ruzhansky.
\newblock {F}ourier multipliers, symbols and nuclearity on compact manifolds.
\newblock {\em {J}. {A}nal. {M}ath.}, 135(2):757$–$800, 2018.

\bibitem[DR21]{dr:intsc}
J.~Delgado and M.~Ruzhansky.
\newblock {S}chatten-von {N}eumann classes of integral operators.
\newblock {\em J. Math. Pures Appl.}, 154: 1--29, 2021.

\bibitem[GG18]{gg:qq1}
F.~Genovese and S.~Gogioso.
\newblock {Q}uantum field theory in categorical quantum mechanics. proceedings
  14th international conference on quantum physics and logic.
\newblock {\em Electron. Proc. Theor. Comput. Sci.}, (266):349--366, 2018.

\bibitem[GG19]{gg:qq2}
F.~Genovese and S.~Gogioso.
\newblock {Q}uantum field theory in categorical quantum mechanics. proceedings
  15th international conference on quantum physics and logic.
\newblock {\em Electron. Proc. Theor. Comput. Sci.}, (287):163--177, 2019.

\bibitem[GS18]{mgas:sp}
M.~Gessner and A.~Smerzi.
\newblock {S}tatistical speed of quantum states: {G}eneralized quantum {F}isher
  information and {S}chatten speed.
\newblock {\em {P}hysical {R}eview {A}.}, 97:022109, 2018.

\bibitem[Gui72]{gui:hi}
A.~Guichardet.
\newblock {\em {S}ymmetric {H}ilbert Spaces and Related Topics. {L}ecture notes
  in mathematics}, volume 261.
\newblock Springer-Verlag Berlin Heidelberg, 1972.%

\bibitem[Gui66]{gui:1} A. Guichardet.  {\em  Produits tensoriels infinis et representations des relations d'anticommutation.} (French) Ann. Sci. Ecole Norm. Sup. (3) 83, 1–-52, 1966.%

\bibitem[LYLW20]{schit:cc1}
J.~Liu, H.~Yuan, X.-M. Lu, and X.~Wang.
\newblock {Q}uantum {F}isher information matrix and multiparameter estimation.
\newblock {\em J. Phys. A.}, 53(2), 2020.

\bibitem[MuN36]{Murray} F. J. Murray and J. von Neumann. On rings of operators. {\em Ann. of Math.} (2) 37 (1936), no. 1, 116–229.


\bibitem[Par92]{Pah:book}
R.~K. Parthasarathy.
\newblock {\em {A}n introduction to quantum stochastic calculus. {M}onographs
  in {M}athematics}, volume~85.
\newblock Birk\"auser Verlag, Basel, 1992.

\bibitem[Piet15]{Piet:trace2}
Pietsch, A. Traces and Residues of Pseudo-Differential Operators on the Torus. Integr. Equ. Oper. Theory 83, 1–23, 2015.

\bibitem[Pru81]{prug:b1}
E.~Prugovecki.
\newblock {\em {Q}uantum mechanics in {H}ilbert Spaces: second edition}.
\newblock Academic Press, New York, 1981.

\bibitem[RS80]{r-s:vol1}
M.~Reed and B.~Simon.
\newblock {\em Methods of modern mathematical physics. {I}}.
\newblock Academic Press, Inc. [Harcourt Brace Jovanovich, Publishers], New
  York, second edition, 1980.
\newblock Functional analysis.

\bibitem[RT10a]{rt:book}
M.~Ruzhansky and V.~Turunen.
\newblock {\em Pseudo-differential operators and symmetries. Background
  analysis and advanced topics}, volume~2 of {\em Pseudo-Differential
  Operators. Theory and Applications}.
\newblock Birkh{\"a}user Verlag, Basel, 2010.


\bibitem[SLH2024]{pnetw:ko} Jin Young Shin, Jae Yong Lee and Hyung Ju Hwang. Pseudo-Differential Neural Operator: Generalized Fourier
{N}eural Operator for Learning Solution Operators of Partial
Differential Equations. {\em Transactions on Machine Learning Research}, 3, 2024


\bibitem[Sch43]{Sch}
  R. Schatten. On the direct product of Banach spaces. {\em Trans. Amer. Math. Soc.} 53 (1943), 195–217

\bibitem[Sch46]{Sch1}
  R. Schatten. The cross-space of linear transformations. {\em Ann. of Math.} (2) 47 (1946), 73–84.

   
\bibitem[SchNe46]{SN}
    R. Schatten and J. von Neumann. The cross-space of linear transformations. II. {\em Ann. of Math.} (2) 47 (1946), 608–630. 

\bibitem[Sch70]{Sch2}
 R. Schatten. Norm ideals of completely continuous operators. Second printing. Ergebnisse der Mathematik und ihrer Grenzgebiete, Band 27 Springer-Verlag, Berlin-New York 1970 vii+81 pp.

\bibitem[Ne39]{Ne}
 J. von Neumann. On infinite direct products. {\em Compositio Math}. 6 (1939), 1–77.

\bibitem[Wea01]{weav:book}
N.~Weaver.
\newblock {\em {M}athematical Quantization}.
\newblock Chapman \& Hall/ Crc, 2001.


\bibitem[Wei80]{weid:b1}
J.~Weidmann.
\newblock {\em {L}inear Operators in Hilbert Spaces}, volume~68.
\newblock Springer-Verlag, New York, 1980.

\end{thebibliography}


\end{document}
**
\begin{rem} It is well known that infinite tensors of bounded operators are bounded. However, this property is not longer valid for compact operators. Consider for instance the sequence of Hilbert spaces $H_j=\ce$, and the operators $A_j=I_j$, the identity operator on  $\ce$. 
    
\end{rem}


****


That  $A_m^p$ is a trace class operator for all $1<p<\infty$, follows from the fact that $A_m$ is a positive definite trace class operator, since its symbol $a_m(x)\beta_m(j)$ begets a composition of a multiplication operator that is bounded with a Fourier multiplier that is trace class. Indeed, since $a_m\in L^{\infty}(\T)$,then the multiplication operator with symbol $a_m$ begets a bounded operator on $L^{2}(\T)$. On the other hand, since $0\leq\beta\in \ell^1(\zet)$, then $\beta(D)\in S_1(L^{2}(\T)$. It also follows that    $A_m^p$ is a trace class operator for all $1<p<\infty$, since the symbol of $A_m^p$

\in L^{\infty}$\\ 


Since $\T$ is of finite measure we have $L^{\infty}(\T)\subset L^{p}(\T)\subset L^{1}(\T)$ for all $1<p<\infty$. On the other hand, $\ell^{1}(\zet)\subset \ell^{p}(\zet)\subset \ell^{\infty}(\zet) $ for all $1<p<\infty$. Thus, when taking the limit when $p\rightarrow 1$, we see that 





Since in the above Theorem one is only interested in $p>1$ for $p$ close to 1. It is not restricted to consider Pietsch definition  








Countably tensors of operators have a wide variety of applications, in particular as an application of the trace formula \eqref{trinf2} for a countably family of trace class operators, we will expand  the {\em  Riemann zeta function} $\zeta(z)$ as the  trace of the tensor of a family of this kind.\\

\begin{thm} Let $\Hi$ be  a complex separable Hilbert space and  $\{e_j\}_{j=1}^{\infty}$ an orthonormal basis of $\Hi$. We denote by $\mathbb{P}$ the set of prime numbers, and we enumerate them in increasing order so that $\mathbb{P}=\{p_1,p_2,\dots\}$. For each $j\in\{1,2,\dots\}$ and $z\in\ce$, such that $\Re{z}>1$, we define the operator $\rho_{j,z}$ on $\Hi$ by 
\[\rho_{j,z}:=\sum\limits_{i=1}^\infty(p_j^{-z})^{i-1}|e_i\rangle\langle e_i|.\]
Then, the operator $\bigotimes \limits_{j=1}^\infty\rho_{j,z}$ is well-defined on $\bigotimes \limits_{j=1}^\infty \Hi$ , it is a trace class operator  and
\[\Tr\left(\bigotimes \limits_{j=1}^\infty \rho_{j,z}\right)=\zeta(z).\]
Consequently, the function $t:\{z\in\ce : \Re{z}>1\}\rightarrow \ce$ defined by 
\[t(z)=\Tr\left(\bigotimes \limits_{j=1}^\infty \rho_{j,z}\right)\]
is analytic and it can be analytically extended to $\ce\setminus\{1\}$.
\end{thm}
\begin{proof} Let $\Re{z}>1$, we  first observe that $\rho_{j,z}$ is  a trace class operator and as a consequence of the geometric series formula we get 
\[\Tr(\rho_{j,z})=\sum\limits_{i=1}^\infty(p_j^{-z})^{i-1}=\frac{1}{1-p_j^{-z}}.\]
Therefore, by Corollary \ref{inftetr1} for the case $p=1$,  we have that $\bigotimes \limits_{j=1}^\infty\rho_{j,z}$ is a trace class operator,  and from the Euler product  formula for the Riemann zeta function we obtain 

\begin{align*}\Tr\left(\bigotimes \limits_{j=1}^\infty \rho_{j,z}\right)=&\prod\limits_{j=1}^{\infty}\Tr(\rho_{j,z})\\
=&\prod\limits_{j=1}^{\infty}\frac{1}{1-p_j^{-z}}\\
=&\zeta(z).
\end{align*}
The properties of the function $t$ are inherited from the corresponding ones  of the Riemann zeta function, since they coincide in $\{z\in\ce : \Re{z}>1\}$.
\end{proof}
Let us denote by $\zeta_T$ the {\em zeta function }  associated to an operator $T$, and defined by 
$\zeta_T(z)=\sum\limits_{j=1}^{\infty}\lambda_j^{-z}$, where $\lambda_j$ are the eigenvalues of $T$ in increasing order. We also consider the harmonic oscillator  
$$\mathcal{H}_2:=-\frac{d^2}{dx^2}+x^2,$$
and write $\zeta_2$ instead $\zeta_{\mathcal{H}_2}$.
\begin{thm} 
    
\end{thm}
\begin{proof} \[\zeta_2(s)=(1-2^{-s})\zeta(s).\]
    
\end{proof}






\subsection{Fock spaces}

We have reviewed the concept of countable Hilbert tensor products in the previous section which is associated to describe systems of infinitely many particles. We now recall a fundamental class of spaces in Quantum Mechanics, the so called {\em Fock spaces } which are  useful as typical state spaces for gases of particles, thermal bath, etc (cf. \cite{},  \cite{}).  These spaces have been built in order to represent the state space for a system of 
 of indefinite and variable number of identical particles as an electron gas, photons, among others. \\

The basic physical idea hidden behind their definition is the following.
If $\Hi$ is the Hilbert space describing the state space for one particle, then
$\Hi\otimes \Hi$ describes the state space for two particles of the same type. The
cspace $\Hi^{\otimes^n}=\bigotimes_{j=1}^n \Hi=\Hi\otimes\cdots\otimes\Hi$, the $n$-fold  tensor product, describes the state space for
$n$ such particles. Finally the space 
\beq\Gamma_f(\Hi):=\bigoplus_{n\in \ene}\Hi^{\otimes^n}\label{ffock}\eq
describes a system where there  
can be any number of such particles which can disappear (be annihilated)
or appear (be created). The space \eqref{ffock} is called 
the {\em full or free Fock space of } $\Hi$. 
But depending on the type of particles (bosons or
fermions), there are some symmetries which force to look at certain subspaces
of the {\em full Fock space } $\Gamma_f(\Hi)$. We now briefly describe those ones. \\

We start with the symmetric one, which arises as the state space for  {\em bosonic } systems. The {\em symmetric Fock space over } $\Hi$ or the {\em Boson  Fock space over } $\Hi$ is defined as follows. For any $n\geq 1$ we consider $\Hi^{\otimes^n}$. For $u_1,\dots, u_n\in\Hi$ we define the {\em n-fold symmetric tensor product  } by 
\beq u_1\circ\cdots\circ u_n=\frac{1}{n!}\sum_{\sigma\in \mathcal{S}_n}u_{\sigma(1)}\otimes\cdots\otimes u_{\sigma(n)}\, ,\eq
where $\mathcal{S}_n$ is the group of permutations of $\{1,2,\dots, n\}.$\\

The closed subspace of  $\Hi^{\otimes^n}$ generated by the $ u_1\circ\cdots\circ u_n$ is denoted by  $\Hi^{on}$ and is called the  {\em n-fold symmetric tensor product of } $\Hi$. \\



For the proof see \cite{weav:book}.***More details ..for this and this Proposition***\\


The {\em symmetric (or Boson) Fock space over } $\Hi$ is the space
 \beq\Gamma_s(\Hi):=\bigoplus_{n\in \ene}\Hi^{\circ n}.\label{bfock}\eq


On the other hand, the {\em antisymmetric tensor product }
 is defined by
\beq u_1\wedge\cdots\wedge u_n=\frac{1}{n!}\sum_{\sigma\in \mathcal{S}_n}\epsilon_{\sigma}u_{\sigma(1)}\otimes\cdots\otimes u_{\sigma(n)}\, ,\eq
where $\epsilon_{\sigma}$ is the signature of the permutation $\sigma$.\\

The closed subspace of  $\Hi^{\otimes^n}$ generated by the $ u_1\wedge\cdots\wedge u_n$
 is denoted by  $\Hi^{\wedge n}$ and is called  the {\em n-fold antisymmetric tensor product of } $\Hi$.\\


Again, 



See (cf. \cite{weav:book}).\\

The {\em antisymmetric (or Fermion) Fock space over } $\Hi$ is the space
 \beq\Gamma_a(\Hi):=\bigoplus_{n\in \ene}\Hi^{\wedge n}.\label{afock}\eq

 
As a convention for $n=0$ we put 
 \[\Hi^{\otimes 0}=\Hi^{\circ 0}=\Hi^{\wedge 0}=\ce.\] 
In physics the spaces $\Hi^{\otimes n}, \Hi^{\circ n}, \Hi^{\wedge n}$ are called the
$n$-{\em particle spaces} and separable spaces as well $\Gamma_a(\Hi)$ and $\Gamma_s(\Hi)$. \\

We point out that the element $1\in\ce$ plays an important role when seen as an element of a Fock space. We denote it by $\Omega$ and call it the {\em vacuum vector.}\\

We now equip $\Hi^{\wedge n}, \Hi^{\circ n}$ with suitable scalar products.  
The scalar product 
\[\langle u_1\wedge\cdots \wedge u_n, v_1\wedge\cdots v_n\rangle =\frac{1}{(n!)^2}\sum\limits_{\sigma, \tau\in S_n}\epsilon_{\sigma}\epsilon_{\tau}\langle u_{\sigma(1)},v_{\tau(1)} , \rangle\cdots\langle u_{\sigma(n)},v_{\tau(n)}  \rangle\]
which by the properties of determinants is equals to
\[\frac{1}{n!}\Det(\langle u_i,v_j\rangle)_{ij}.\]
The factor $n!$ can be removed by endowing $\Hi^{\wedge n}$ with a different scalar product from the one induced by $\Hi^{\otimes n}$.  So we define:\\
\beq \langle u_1\wedge\cdots \wedge u_n, v_1\wedge\cdots v_n\rangle_{\wedge}:= \Det(\langle u_i,v_j\rangle)_{ij}.\eq
Thus we obtain
\beq\| u_1\wedge\cdots \wedge u_n\|_{\wedge}^2=n!\| u_1\wedge\cdots \wedge u_n\|_{\otimes}^2.\eq
On the other hand, for the space $\Hi^{\circ n}$ we define
 \beq \langle u_1\circ\cdots \circ u_n, v_1\circ\cdots \circ v_n\rangle_{\circ}:= Per(\langle u_i,v_j\rangle)_{ij},\eq
where $Per$ denotes the {\em permanent } of the matrix (that is, the determinant without the minus signs). Hence
\beq\|  u_1\circ\cdots \circ u_n\|_{\circ}^2=n!\| u_1\circ\cdots \circ u_n\|_{\otimes}^2.\eq

Therefore

\begin{prop}
If $\Hi$ be a separable Hilbert space and $\{e_i\}$ be an orthonormal basis of $\Hi$, for $n\geq 1$, we have
$$e_{i_1}\circ\cdots\circ e_{i_n}$$ with $i_1 \leq i_2 \leq \cdots \leq i_n$  in form an orthogonal basis of $\Hi^{on}$.
\end{prop}

\begin{prop}
If $\Hi$ be a separable Hilbert space and $\{e_i\}$ be an orthonormal basis of $\Hi$, for $n\geq 1$, we have
$$e_{i_1}\wedge\cdots\wedge e_{i_n}$$ with $i_1 < i_2 < \cdots < i_n$  in form an orthogonal basis of $\Hi^{\wedge n}$.
\end{prop}

It should be understood that in the definition of  $\Gamma_f(\Hi), \Gamma_s(\Hi), \Gamma_a(\Hi)$ each of the spaces $\Hi^{\otimes n}, \Hi^{\circ n}, \Hi^{\wedge n}$
 is endowed with the respective scalar product $ .$


******************

Let $()$ be an orthonormal basis of $\Hi$\\



If X is a $\sigma$-finite measure space then we regard $\Gamma_s(L^2(X))$ as a measurable tensor product of the Hilbert spaces $l_2(\mathbb{N})$ over the index set $X$, and we regard  $\Gamma_s(L^2(X))$ as a measurable tensor product of the Hilbert spaces $\mathbb{C}^2$ over the index set $X$. 


***************************



Since at the end Fock spaces are determined by a single state Hilbert space $\Hi$, 
 we now give an important example of such space. 
 
 
 
 
 
 
 
 If $\Hi==L^2(\ar)$ then,
 
 
 
 
***************************
 
 In the case of a single Schr\"odinger particle of spin one-half the corresponding Hilbert space is $ \Hi=L^2(\ar^3,\ce^2)$ which is isomorphic to $ \Hi=L^2(\ar^3)\otimes\ce^2$***\\

















The creation
and annihilation operators are used to account for the introduction
and removal of particles, allowing us to describe a system with a
variable number of particles.

  

***When

The Fock spaces are used in... symmetric and antisimmetric Fock spaces arise as the states spaces for {\em bosonic } and  {\em fermionic} systems respectively.\\




********





In quantum mechanics identical particles are indistinguishable. Thus in a
Fock space, all particles must be identical. In order to consider multiple types
of particles, one must take the tensor product of different Fock spaces - one for
each species


We ...Fock sspaces associated to compact manifolds and in particular compact Lie groups as ambient spaces****



*****

\newpage

Sean  las permutaciones 
$\epsilon=\left(
 \begin{array}{cc}
  1 & 2\\
  1 & 2\\
  \end{array}
\right)$, $\sigma=\left(
 \begin{array}{cc}
  1 & 2\\
  2 & 1\\
  \end{array}
\right)$,
Entonces dados $u_1,u_2 \in H$, tenemos que 
\[u_1 \circ u_2 = \frac{1}{2}[u_{\epsilon(1)}\otimes u_{\epsilon(2)}+u_{\sigma(1)}\otimes u_{\sigma(2)}]=\frac{1}{2}[u_1\otimes u_2+u_{2}\otimes u_{1}].\]

Para $i=1,2$, tenemos que $u_i= \sum\limits_{m=1}^\infty a_{im}e_m^{(i)}$, as\'i que
\[u_{\epsilon(i)}=\sum\limits_{m=1}^\infty a_{\epsilon(i)m}e_m^{(\epsilon(i))} \quad \textnormal{y} \quad u_{\sigma(i)}=\sum\limits_{m=1}^\infty a_{\sigma(i)m}e_m^{(\sigma(i))}\]

\begin{align*}
   u_{\epsilon(1)}\otimes u_{\epsilon(2)} & = \left(\sum\limits_{m=1}^\infty a_{\epsilon(1)m}e_m^{(\epsilon(1))}\right)\otimes\left(\sum\limits_{m=1}^\infty a_{\epsilon(2)m}e_m^{(\epsilon(2))}\right) \\ 
&= \sum\limits_{m=1}^\infty a_{\epsilon(1)1}a_{\epsilon(2)m}e_1^{(\epsilon(1))}\otimes e_m^{(\epsilon(2))} + \sum\limits_{m=1}^\infty a_{\epsilon(1)2}a_{\epsilon(2)m}e_2^{(\epsilon(1))}\otimes e_m^{(\epsilon(2))} \\
& + \sum\limits_{m=1}^\infty a_{\epsilon(1)3}a_{\epsilon(2)m}e_3^{\epsilon(1)}\otimes e_m^{(\epsilon(1))}+ \ldots 
\end{align*}

Por otro lado 

\begin{align*}
   u_{\sigma(1)}\otimes u_{\sigma(2)} & = \left(\sum\limits_{m=1}^\infty a_{\sigma(1)m}e_m^{(\sigma(1))}\right)\otimes\left(\sum\limits_{m=1}^\infty a_{\sigma(2)m}e_m^{(\sigma(2))}\right) \\ 
&= \sum\limits_{m=1}^\infty a_{\sigma(1)1}a_{\sigma(2)m}e_1^{(\sigma(1))}\otimes e_m^{(\sigma(2))} + \sum\limits_{m=1}^\infty a_{\sigma(1)2}a_{\sigma(2)m}e_2^{(\sigma(1))}\otimes e_m^{(\sigma(2))} \\
& + \sum\limits_{m=1}^\infty a_{\sigma(1)3}a_{\sigma(2)m}e_3^{\sigma(1)}\otimes e_m^{(\sigma(1))}+ \ldots 
\end{align*}
Por tanto 

\begin{align*}
u_1 \circ u_2 & = \frac{1}{2}[u_{\epsilon(1)}\otimes u_{\epsilon(2)}+u_{\sigma(1)}\otimes u_{\sigma(2)}] \\
& = \frac{1}{2} (\sum\limits_{m=1}^\infty a_{\epsilon(1)1}a_{\epsilon(2)m}e_1^{(\epsilon(1))}\otimes e_m^{(\epsilon(2))} + \sum\limits_{m=1}^\infty a_{\sigma(1)1}a_{\sigma(2)m}e_1^{(\sigma(1))}\otimes e_m^{(\sigma(2))} \\
& + \sum\limits_{m=1}^\infty a_{\epsilon(1)2}a_{\epsilon(2)m}e_2^{(\epsilon(1))}\otimes e_m^{(\epsilon(2))} + \sum\limits_{m=1}^\infty a_{\sigma(1)2}a_{\sigma(2)m}e_2^{(\sigma(1))}\otimes e_m^{(\sigma(2))} \\
& + \sum\limits_{m=1}^\infty a_{\epsilon(1)3}a_{\epsilon(2)m}e_3^{\epsilon(1)}\otimes e_m^{(\epsilon(1))} + \sum\limits_{m=1}^\infty a_{\sigma(1)3}a_{\sigma(2)m}e_3^{\sigma(1)}\otimes e_m^{(\sigma(1))}+ \ldots)\\
&= \frac{1}{2}(\sum\limits_{m=1}^\infty [a_{\epsilon(1)1}a_{\epsilon(2)m}+a_{\sigma(1)1}a_{\sigma(2)m}]e_1 \otimes e_m + \sum\limits_{m=1}^\infty [a_{\epsilon(1)2}a_{\epsilon(2)m}+a_{\sigma(1)2}a_{\sigma(2)m}]e_2 \otimes e_m \\
& + \sum\limits_{m=1}^\infty [a_{\epsilon(1)3}a_{\epsilon(2)m}+a_{\sigma(1)3}a_{\sigma(2)m}]e_3 \otimes e_m +\dots)
\end{align*}
Por otro lado 

\begin{enumerate}
\item Si $i=j$, $e_i \circ e_i = \frac{1}{2}(e_i \otimes e_i + e_i \otimes e_i)=e_i \otimes e_i$
\item Para $i< j$, $e_i \circ e_j = \frac{1}{2}(e_i \otimes e_j + e_j \otimes e_i)$

Al desarrollar los sumandos anteriores encontramos elementos de la forma $e_i \otimes e_j$ para $i=j$, $i<j$ e $i>j$. Por lo tanto si desglosamos los términos y reagrupamos nos damos cuenta que $a_{\epsilon(1)j}=a_{\sigma(2)j}$ y que $a_{\epsilon(2)j}=a_{\sigma(1)j}$
\begin{align*}
u_1 \circ u_2 & = \frac{1}{2}[u_{\epsilon(1)}\otimes u_{\epsilon(2)}+u_{\sigma(1)}\otimes u_{\sigma(2)}] \\
& = \frac{1}{2}\sum\limits_{j=1}^{\infty}[a_{\epsilon(1)j}a_{\epsilon(2)j}+a_{\sigma(1)j}a_{\sigma(2)j}]e_j \otimes e_j \\
& + \frac{1}{2}\sum\limits_{i,j, i<j}^{\infty}[a_{\epsilon(1)i}a_{\epsilon(2)j}+a_{\sigma(1)i}a_{\sigma(2)j}]e_i \otimes e_j + \frac{1}{2}\sum\limits_{i,j, i>j}^{\infty} [a_{\epsilon(1)j}a_{\epsilon(2)i}+a_{\sigma(1)j}a_{\sigma(2)i}] e_j \otimes e_i \\
& = \frac{1}{2}\sum\limits_{j=1}^{\infty}[a_{1j}a_{2j}+a_{2j}a_{1j}]e_j \otimes e_j + \frac{1}{2}\sum\limits_{i,j i\neq j}^{\infty}[a_{\epsilon(1)i}a_{\epsilon(2)j}+a_{\sigma(1)i}a_{\sigma(2)j}](e_i \otimes e_j +  e_j \otimes e_i)\\
&=\frac{1}{2}\sum\limits_{j=1}^{\infty}[2a_{1j}a_{2j}]e_j \otimes e_j +\sum\limits_{i,j i\neq j}^{\infty}[a_{\epsilon(1)i}a_{\epsilon(2)j}+a_{\sigma(1)i}a_{\sigma(2)j}]\frac{1}{2}(e_i \otimes e_j +  e_j \otimes e_i)\\
& =  \sum\limits_{j=1}^{\infty}[a_{1j}a_{2j}]e_j \circ e_j + \sum\limits_{i,j i< j}^{\infty} [a_{\epsilon(1)i}a_{\epsilon(2)j}+a_{\sigma(1)i}a_{\sigma(2)j}](e_i \circ e_j)
\end{align*}

Es decir  que el generado por la familia  $$\{e_i \circ e_j, \quad i \leq j, \quad i,j\in \mathbb{N}\}$$ son los tensores sim\'etricos $u_1\circ u_2$.


On the other hand, for the space $\Hi^{\circ n}$ we define
 \beq \langle u_1\circ\cdots \circ u_n, v_1\circ\cdots \circ v_n\rangle_{\circ}:= \frac{1}{n!}Per(\langle u_i,v_j\rangle)_{ij},\eq
where $Per$ denotes the {\em permanent } of the matrix (that is, the determinant without the minus signs). Hence
\beq\|  u_1\circ\cdots \circ u_n\|_{\circ}^2=n!\| u_1\circ\cdots \circ u_n\|_{\otimes}^2.\eq


\begin{align*} 
\langle e_{i_l}\circ e_{j_k}, e_{i_l}\circ e_{j_k}\rangle_{\circ} & :=  \left\langle \frac{1}{2}( e_{i_l}\otimes e_{j_k} + e_{j_k}\otimes e_{i_l}), \frac{1}{2}( e_{i_l}\otimes e_{j_k} + e_{j_k}\otimes e_{i_l})\right \rangle\\
& = \frac{1}{4}\left\langle e_{i_l}\otimes e_{j_k} + e_{j_k}\otimes e_{i_l}, e_{i_l}\otimes e_{j_k} + e_{j_k}\otimes e_{i_l}\right \rangle\\
& = \frac{1}{4}(\|e_{i_l}\otimes e_{j_k} \|^2 + \|e_{j_k}\otimes e_{i_l} \|^2 + \langle e_{i_l}\otimes e_{j_k},e_{j_k}\otimes e_{i_l} \rangle + \langle e_{j_k}\otimes e_{i_l},e_{i_l}\otimes e_{j_k} \rangle )\\
& = \frac{1}{4}(2\| e_{i_l}\|^{2}\| e_{j_k}\|^2+\langle e_{i_l},e_{j_k}\rangle \langle e_{j_k}, e_{i_l} \rangle + \langle e_{j_k}, e_{i_l}\rangle \langle e_{i_l}, e_{j_k}\rangle)\\
& = \frac{1}{4}(2+\langle e_{i_l},e_{j_k}\rangle \langle e_{j_k}, e_{i_l} \rangle + \langle e_{j_k}, e_{i_l}\rangle \langle e_{i_l}, e_{j_k}\rangle)
\end{align*}
En este caso si $i_l < j_k$ tenemos que $\langle e_{i_l}\circ e_{j_k}, e_{i_l}\circ e_{j_k}\rangle_{\circ} =\frac{1}{4}(2)=\frac{1}{2} =  \|e_{i_l}\circ e_{j_k}\|^2_{\circ}$  es  decir que 
 $\|e_{i_l}\circ e_{j_k}\|_{\circ} = \sqrt{2}$ y 
si $i_l = j_k$ tenemos que $\langle e_{i_l}\circ e_{j_k}, e_{i_l}\circ e_{j_k}\rangle_{\circ} = \frac{1}{4}4=1 = \|e_{i_l}\circ e_{j_k}\|^2_{\circ}$.

Supongamos que $i_l \neq s_m\leq j_k$.

\begin{align*} 
\langle e_{i_l}\circ e_{j_k}, e_{s_m}\circ e_{j_k}\rangle_{\circ} & := \left\langle \frac{1}{2}( e_{i_l}\otimes e_{j_k} + e_{j_k}\otimes e_{i_l}), \frac{1}{2}( e_{s_m}\otimes e_{j_k} + e_{j_k}\otimes e_{s_m})\right \rangle\\
& = \frac{1}{4}\left\langle e_{i_l}\otimes e_{j_k} + e_{j_k}\otimes e_{i_l}, e_{s_m}\otimes e_{j_k} + e_{j_k}\otimes e_{s_m}\right \rangle\\
& = \frac{1}{4}( \langle e_{i_l}\otimes e_{j_k},e_{s_m}\otimes e_{j_k} \rangle +\langle e_{i_l}\otimes e_{j_k},e_{j_k}\otimes e_{s_m} \rangle + \langle e_{j_k}\otimes e_{i_l},e_{s_m}\otimes e_{j_k} \rangle\\
& + \langle e_{j_k}\otimes e_{i_l},e_{j_k}\otimes e_{s_m} \rangle )\\
& = \frac{1}{4}(\langle e_{i_l},e_{s_m}\rangle\langle e_{j_k},e_{j_k}\rangle +\langle e_{i_l},e_{j_k}\rangle\langle e_{j_k},e_{s_m}\rangle +\langle e_{j_k},e_{s_m}\rangle \langle e_{i_l}, e_{j_k} \rangle + \langle e_{j_k}, e_{j_k}\rangle \langle e_{i_l}, e_{s_m}\rangle)\\ 
& = 0
\end{align*}

Ahora consideremos que 
 
\[\langle e_{j_1}\circ e_{j_2}\circ \ldots  \circ e_{j_n}, e_{j_1}\circ e_{j_2}\circ \ldots \circ e_{j_n}\rangle \] \[= \left\langle \frac{1}{n!}\sum \limits_{s=1}^{n!}e_{j_{\epsilon_s(1)}}\otimes  e_{j_{\epsilon_s(2)}}\otimes \ldots \otimes e_{j_{\epsilon_s(n)}}, \frac{1}{n!}\sum \limits_{s=1}^{n!}e_{j_{\epsilon_s(1)}}\otimes  e_{j_{\epsilon_s(2)}}\circ \ldots \otimes e_{j_{\epsilon_s(n)}} \right \rangle \] \[= \frac{1}{n!^2}\left\langle\sum \limits_{s=1}^{n!}e_{j_{\epsilon_s(1)}}\otimes  e_{j_{\epsilon_s(2)}}\circ \ldots \otimes e_{j_{\epsilon_s(n)}}, \sum \limits_{s=1}^{n!}e_{j_{\epsilon_s(1)}}\otimes  e_{j_{\epsilon_s(2)}}\otimes \ldots \otimes e_{j_{\epsilon_s(n)}} \right \rangle\] \[=\frac{1}{n!^2}\left\langle\sum \limits_{s=1}^{n!}e_{j_{\epsilon_s(1)}}\otimes  e_{j_{\epsilon_s(2)}}\otimes \ldots \otimes e_{j_{\epsilon_s(n)}}, e_{j_{\epsilon_1(1)}}\otimes  e_{j_{\epsilon_1(2)}}\otimes \ldots \otimes e_{j_{\epsilon_1(n)}} \right \rangle +\] \[\frac{1}{n!^2}\left\langle\sum \limits_{s=1}^{n!}e_{j_{\epsilon_s(1)}}\otimes  e_{j_{\epsilon_s(2)}}\otimes \ldots \otimes e_{j_{\epsilon_s(n)}}, e_{j_{\epsilon_2(1)}}\otimes  e_{j_{\epsilon_2(2)}}\otimes \ldots \otimes e_{j_{\epsilon_2(n)}} \right \rangle +\] \[\ldots + \frac{1}{n!^2}\left\langle\sum \limits_{s=1}^{n!}e_{j_{\epsilon_s(1)}}\otimes  e_{j_{\epsilon_s(2)}}\otimes \ldots \otimes e_{j_{\epsilon_{s}(n)}}, e_{j_{\epsilon_{n!}(1)}}\otimes  e_{j_{\epsilon_{n!}(2)}}\otimes \ldots \otimes e_{j_{\epsilon_{n!}(n)}} \right \rangle \] \[= \frac{1}{n!^2}\sum\limits_{s=1}^{n!}\langle e_{j_{\epsilon_s(1)}},e_{j_{\epsilon_{1}(1)}}\rangle \langle e_{j_{\epsilon_s(2)}}, e_{j_{\epsilon_{1}(2)}}\rangle \ldots \langle e_{j_{\epsilon_{s}(n)},e_{j_{\epsilon_{1}(n)}}}\rangle +\] \[\frac{1}{n!^2}\sum\limits_{s=1}^{n!}\langle e_{j_{\epsilon_s(1)}},e_{j_{\epsilon_{2}(1)}}\rangle \langle e_{j_{\epsilon_s(2)}}, e_{j_{\epsilon_{2}(2)}}\rangle \ldots \langle e_{j_{\epsilon_{s}(n)},e_{j_{\epsilon_{2}(n)}}}\rangle + \] \[\ldots +\frac{1}{n!^2}\sum\limits_{s=1}^{n!}\langle e_{j_{\epsilon_s(1)}},e_{j_{\epsilon_{n!}(1)}}\rangle \langle e_{j_{\epsilon_s(2)}}, e_{j_{\epsilon_{n!}(2)}}\rangle \ldots \langle e_{j_{\epsilon_{s}(n)},e_{j_{\epsilon_{n!}(n)}}}\rangle   
\]

Para ello hay que considerar dos casos \\
1. Si $j_1=j_2 = \ldots = j_n$, entonces $e_{j_{\epsilon_s(n)}}=e_{j_n}$ para todo $n$. Por tanto toda la suma sobrevive y habr\'ia $n!$ t\'erminos.
\[\langle e_{j_1}\circ e_{j_2}\circ \ldots  \circ e_{j_n}, e_{j_1}\circ e_{j_2}\circ \ldots \circ e_{j_n}\rangle = \frac{1}{n!^2}(n!+n!+n!+\ldots +n!)\] ($n!$ t\'erminos)  Entonces \[\langle e_{j_1}\circ e_{j_2}\circ \ldots  \circ e_{j_n}, e_{j_1}\circ e_{j_2}\circ \ldots \circ e_{j_n}\rangle= \frac{1}{n!^2}(n!^2)= 1=\|e_{j_1}\circ e_{j_2}\circ \ldots  \circ e_{j_n}\|^2\]
2. Si $j_1< j_2 < \ldots < j_n$ entonces los productos punto sobreviven solo cuando $s=1, s=2, \ldots s=n!$. Por tanto  \[\langle e_{j_1}\circ e_{j_2}\circ \ldots  \circ e_{j_n}, e_{j_1}\circ e_{j_2}\circ \ldots \circ e_{j_n}\rangle = \frac{1}{n!^2}(n!)= \frac{1}{n!}= \|e_{j_1}\circ e_{j_2}\circ \ldots  \circ e_{j_n}\|^2\]
As\'i que \[\|e_{j_1}\circ e_{j_2}\circ \ldots \circ e_{j_n}\|=\frac{1}{\sqrt{n!}}\]
\end{document}




Sean $f\in L^2(X, \mathcal{A},\mu)$ y $g\in L^2(Y, \mathcal{B},\nu)$. Entonces $$T: L^2(X, \mathcal{A},\mu) \otimes L^2(Y, \mathcal{B},\nu) \to L^2(X \times Y, \mathcal{A}\otimes \mathcal{B},\mu\otimes \nu)$$ definido por $T(f\otimes g)(x,y)=f(x)g(y)$ es un isomorfismo y un operador de Hilbert Schmid cuyo n\'ucleo $k(x,y)=f(x)g(y)$.

En efecto,
Sea $u=\sum\limits_{i=1}^k f_i\otimes g_i$ y mostremos que $$\|T(u)\|=\|u\|.$$
\begin{align*}
\|T(u)\|^2 & = \langle T(u) , \overline{T(u)}\rangle\\
&= \int_{X\times Y}T(u)\overline{T(u)}d(\mu\otimes \nu)(x,y)\\
& = \int_{X\times Y} \sum\limits_{i=1}^k f_i(x)g_i(y)\overline{\sum\limits_{i=1}^k f_i(x)g_i(y)}d(\mu\otimes \nu)(x,y)\\
& = \int_{X\times Y} \sum\limits_{i=1}^k f_i(x)g_i(y)\sum\limits_{i=1}^k \overline{f_i(x)g_i(y)}d(\mu\otimes \nu)(x,y)\\
& = \sum\limits_{i=1}^k\sum\limits_{j=1}^k\int_{X\times Y}f_i(x)\overline{f_j(x)}g_i(y)\overline{g_j(y)}d(\mu\otimes \nu)(x,y)\\
& =\sum\limits_{i=1}^k\sum\limits_{j=1}^k\left(\int_{X}f_i(x)\overline{f_j(x)}d\mu(x)\int_{Y}g_i(y)\overline{g_j(y)}d\nu(y)\right)\\
& = \sum\limits_{i=1}^k\sum\limits_{j=1}^k\langle f_i , f_j\rangle\langle g_i , g_j\rangle
\end{align*}

Luego $$\|T(u)\|= \left(\sum\limits_{i=1}^k\sum\limits_{j=1}^k\langle f_i , f_j\rangle\langle g_i , g_j\rangle\right)^ {1/2}= \|u\|. $$

Si tomamos a $X=Y=\mathbb{N}$ y $\mu,\nu$ como las medidas de contar, entonces

\begin{align*}
\|T(u)\|^2 & = \sum\limits_{i=1}^k\sum\limits_{j=1}^k\left(\int_{\mathbb{N}}f_i(x)\overline{f_j(n)}d\mu(n)\int_{\mathbb{N}}g_i(m)\overline{g_j(m)}d\nu(m)\right)\\
& = \sum\limits_{i=1}^k\sum\limits_{j=1}^k \left( \sum\limits_{n\in \mathbb{N}}f_i(x)\overline{f_j(n)}\sum\limits_{m\in \mathbb{N}}g_i(m)\overline{g_j(m)}\right)\\
& = \sum\limits_{i=1}^k\sum\limits_{j=1}^k\langle f_i , f_j\rangle\langle g_i , g_j\rangle 
\end{align*}













\end{document}
M. Chatzakou, J. Delgado and M. Ruzhansky. On a class of anharmonic oscillators II: General case.  Bulletin des Sciences Mathématiques. 180,  art.no. 103196. November, 2022


 \begin{rem} Let $\Hi$ be a complex separable Hilebrt space, and let $\{e_n:n\in \ene\}$ be an orthonormal basis of $\Hi$.  Consider the sequence of positive compact operators $A_N$ defined on $\Hi$ determined by the corresponding sequence of eigenvalues $\{\frac{N}{n}\}_{n}$, i.e.,  $A_N(e_n)=\frac{N}{n}e_n$. Then, given any stalilising sequence ....it is not possible to define  the value of ...on the basis....as a natural extension of . It becomes neccesary to impose the following condition on the sequence of operators ***** \\

     For diagonalizable operators, and assuming some type of behaviour for the norms of the sequence of operators we can define the tensor of...***
 \end{rem}



 On the other hand,
\begin{align*} \prod\limits_{j=1}^{n}\|A_j\|_{\mathcal{L}(\Hi_j)}& \leq ****\\
& \leq \|\bigotimes \limits_{j=1}^{\infty}A_j\|_{\mathcal{L}\left(\bigotimes \limits_{j=1}^{\infty}\Hi_j\right)}.
  \end{align*}
  




\begin{ex} There is a bulk of well-known formulas on infinite products relating to special functions. For instance, it is known that  
 \[\prod\limits_{n=1}^{\infty}(1+\frac{1}{n^{p}}),\]
 converges if $p>1.$\\
Let us consider  a positive real number $p$ with $p>1$. For each $k\in\{1,2,\dots\}$, we define   a bounded diagonal operator on $\Hi_k$ by 
\[A_{k}(e_j^{k}):=(1+\frac{1}{jk^p})e_j^{k},\,\, \mbox{ for }j=1,2\dots.\]
 Then, $\|A_{k}\|_{ \mathcal{L}(\Hi_k)}=1+\frac{1}{k^p}$   and
 \[\|\bigotimes \limits_{k=1}^{\infty}A_{k}\|_{\mathcal{L}(\Hi_1\otimes\cdots\otimes \Hi_n\otimes\cdots)}=\prod\limits_{k=1}^{\infty}\|A_{k}\|_{\mathcal{L}(\Hi_k)}=\prod\limits_{k=1}^{n}(1+\frac{1}{k^p})=\frac{\Gamma(n+z+1)}{n!\Gamma(z+1)}.\]



The following example ..Let $z>0$, for each $k\in\{1,2,\dots,n\}$, we define  $A_{k,z}\in \mathcal{L}(\Hi_k)$ as a bounded diagonal operator by 
\[A_{k,z}(e_j^{k}):=(1+\frac{z}{kj})e_j^{k},\,\, \mbox{ for }j=1,2\dots.\]
 Then, $\|A_{k,z}\|=1+\frac{z}{k}$   and
 \[\|\bigotimes \limits_{k=1}^nA_{k,z}\|_{\mathcal{L}(\Hi_1\otimes\cdots\otimes \Hi_n)}=\prod\limits_{k=1}^{n}\|A_{k,z}\|_{\mathcal{L}(\Hi_k)}=\prod\limits_{k=1}^{n}(1+\frac{z}{k})=\frac{\Gamma(n+z+1)}{n!\Gamma(z+1)}.\]
 For the last identity we have used ****
\end{ex}


The following condition will be useful for the study of compactness of infinite tensors.
\begin{defn} \label{matrixef1} Let $\mathcal{M}=(\mathcal{M}_{k\,j})_{(j,k)\in \ene\times\ene}$ be an infinite matrix of complex entries. \\

\noindent $(\Pi_0)$\,\, We say that $\mathcal{M}$ satisfies the condition $(\Pi_0)$ if 
\[\lim_{n\rightarrow\infty}\mathcal{M}_{k_1\,1}\mathcal{M}_{k_1\,2}\cdots\mathcal{M}_{k_n\,n}=0,\]
for any sequence $(k_j)_{j}$.
\end{defn}%
\begin{defn}
For a sequence of bounded diagonal operators $A_j$ as in Theorem .. (i), we consider the decompositions \eqref{} for each $A_j$, and define the \em{matrix of eigenvalues} $\mathcal{M}_E$ by:
\begin{equation}\label{ME}
 (\mathcal{M}_E)_{jk}=(\lambda_{k}^{j})_{jk}.
\end{equation}
\end{defn}



Countably tensors of operators have a wide variety of applications, in particular as an application of the trace formula \eqref{trinf2} for a countably family of trace class operators, we will expand  the {\em  Riemann zeta function} $\zeta(z)$ as the  trace of the tensor of a family of this kind.\\

\begin{thm} Let $\Hi$ be  a complex separable Hilbert space and  $\{e_j\}_{j=1}^{\infty}$ an orthonormal basis of $\Hi$. We denote by $\mathbb{P}$ the set of prime numbers, and we enumerate them in increasing order so that $\mathbb{P}=\{p_1,p_2,\dots\}$. For each $j\in\{1,2,\dots\}$ and $z\in\ce$, such that $\Re{z}>1$, we define the operator $\rho_{j,z}$ on $\Hi$ by 
\[\rho_{j,z}:=\sum\limits_{i=1}^\infty(p_j^{-z})^{i-1}|e_i\rangle\langle e_i|.\]
Then, the operator $\bigotimes \limits_{j=1}^\infty\rho_{j,z}$ is well-defined on $\bigotimes \limits_{j=1}^\infty \Hi$ , it is a trace class operator  and
\[\Tr\left(\bigotimes \limits_{j=1}^\infty \rho_{j,z}\right)=\zeta(z).\]
Consequently, the function $t:\{z\in\ce : \Re{z}>1\}\rightarrow \ce$ defined by 
\[t(z)=\Tr\left(\bigotimes \limits_{j=1}^\infty \rho_{j,z}\right)\]
is analytic and it can be analytically extended to $\ce\setminus\{1\}$.
\end{thm}
\begin{proof} Let $\Re{z}>1$, we  first observe that $\rho_{j,z}$ is  a trace class operator and as a consequence of the geometric series formula we get 
\[\Tr(\rho_{j,z})=\sum\limits_{i=1}^\infty(p_j^{-z})^{i-1}=\frac{1}{1-p_j^{-z}}.\]
Therefore, by Corollary \ref{inftetr1} for the case $p=1$,  we have that $\bigotimes \limits_{j=1}^\infty\rho_{j,z}$ is a trace class operator,  and from the Euler product  formula for the Riemann zeta function we obtain 

\begin{align*}\Tr\left(\bigotimes \limits_{j=1}^\infty \rho_{j,z}\right)=&\prod\limits_{j=1}^{\infty}\Tr(\rho_{j,z})\\
=&\prod\limits_{j=1}^{\infty}\frac{1}{1-p_j^{-z}}\\
=&\zeta(z).
\end{align*}
The properties of the function $t$ are inherited from the corresponding ones  of the Riemann zeta function, since they coincide in $\{z\in\ce : \Re{z}>1\}$.
\end{proof}
Let us denote by $\zeta_T$ the {\em zeta function }  associated to an operator $T$, and defined by 
$\zeta_T(z)=\sum\limits_{j=1}^{\infty}\lambda_j^{-z}$, where $\lambda_j$ are the eigenvalues of $T$ in increasing order. We also consider the harmonic oscillator  
$$\mathcal{H}_2:=-\frac{d^2}{dx^2}+x^2,$$
and write $\zeta_2$ instead $\zeta_{\mathcal{H}_2}$.
\begin{thm} 
    
\end{thm}
\begin{proof} \[\zeta_2(s)=(1-2^{-s})\zeta(s).\]
    
\end{proof}

$$\zeta_2(z)$ the zeta function $\zeta_2(z)$ of the  quantum harmonic oscillator, thanks to the identity**** \\

\[\]

we also get ***\\

In the case of compact Lie groups****\\

Let us consider a family of pseudo-differential operators $\{A_j\}$ ...\\

In particular, for $p=1$ the trace class we additionally have 
\[\Tr\left(\bigotimes \limits_{j=1}^\infty A_j\right)=\prod\limits_{j=1}^{\infty}\int\limits_{G}\sum\limits_{\xi}\sigma_j(x,\xi)dx.\]	




******



playground.arduino.cc/Code/LcdBarGraph/




In the case of countably families of Schatten-von Neumann invariant operators, as a consequence of Theorem \ref{inftetr1},   the Corollary above and Theorem 2.5 of \cite{fjdmr:foum} we have :
\begin{cor} \label{invSchatten3} Let  $0<p<\infty$ and   $\Hi_j$  complex separable Hilbert spaces for $j=1, 2, \dots \,$. If $A_j\in S_p(\Hi_j)$ is invariant  for $j=1, 2, \dots \,$; and  ***$\prod\limits_{j=1}^{\infty}\|A_j\|_{S_p(\Hi_j)}\rightarrow 1$. Then  **********
\[\bigotimes \limits_{j=1}^{\infty}A_j\in S_p\left(\bigotimes \limits_{j=1}^{\infty}\Hi_j\right).\]
Moreover
\[\|\bigotimes \limits_{j=1}^{\infty}A_j\|_{S_p\left(\bigotimes \limits_{j=1}^{\infty}\Hi_j\right)}=\prod\limits_{j=1}^{\infty}\|A_j\|_{S_p(\Hi_j)}=\prod\limits_{j=1}^{\infty}\left(\sum\limits_{\ell=1}^{\infty}\|\sigma_j(\ell)\|_{S_p(\Hi_{\ell})}^p \right)^{\frac{1}{p}}.\]
In the case of the  trace class ($p=1$) we additionally have 
\begin{equation}
\Tr\left(\bigotimes \limits_{j=1}^\infty A_j\right)=\prod\limits_{j=1}^{\infty}\Tr(A_j)=\prod\limits_{j=1}^{\infty}\left(\sum\limits_{\ell=1}^{\infty}\Tr(\sigma_j(\ell) )\right).\label{trinf2h}\end{equation}	 
\end{cor}
%Then, the Hilbert space $\Hi_1\otimes\Hi_2$ can be written in the form 
%\[\Hi_1\otimes\Hi_2=\bigoplus\limits_{j,\,k=1}^{\infty}(\Hi_{1,j}\otimes\Hi_{2,k})\]

%**** \[\sigma(j,k):\Hi_{1,j}\otimes\Hi_{2,k}\rightarrow \Hi_{1,j}\otimes\Hi_{2,k}\]









%consider two complex separable Hilbert spaces $\Hi_1, \Hi_2$ decomposed as direct sums in the form
%\[\Hi_1=\bigoplus\limits_{j=1}^{\infty}\Hi_{1,j},\,\, \Hi_2=\bigoplus\limits_{k=1}^{\infty}\Hi_{2,k} \]

%******We dont need invariant here

In particular, for $p=1$ the trace class we additionally have 
\[\Tr\left(\bigotimes \limits_{j=1}^\infty A_j\right)=\prod\limits_{j=1}^{\infty}\int\limits_{G}\sum\limits_{\xi}\sigma_j(x,\xi)dx.\]



We now establish  conditions for the compactness,  membership to Schatten-von Neumann  classes of infinite tensors, the corresponding norms and trace formulae. 
\begin{thm} \label{inftetr1b} Let  $\Hi_j$  be a complex separable Hilbert space for $j\in\{1, 2, \dots \,\}$. We asume that  $A_j\in S_{\infty}(\Hi_j)$  has a diagonal form as in Definition \ref{diagdef1}, for every $j$.\\ 

\noindent {\mbox{(i)}} If the operators $A_j$ are compact diagonal and the  net of partial products $\{\prod\limits_{j\in F}\|A_j\|_{\mathcal{L}(\Hi_j)}\}_{F\in\fou}$  diverges to zero. 
 Then, the operator $\bigotimes \limits_{j=1}^{\infty}A_j$ is compact diagonal.\\ 

\noindent {\mbox{(ii)}} Let  $0<p<\infty$. If $A_j\in S_p(\Hi_j)$ for $j=1, 2, \dots \,$; and assume that the  net of partial products $\{\prod\limits_{j\in F}\|A_j\|_{S_p(\Hi_j)}\}_{F\in\fou}$  converges to zero. 
 Then $\bigotimes \limits_{j=1}^{\infty}A_j$ belongs to $ S_p\left(\bigotimes \limits_{j=1}^{\infty}\Hi_j\right)$.***\\
 
Moreover
\[\|\bigotimes \limits_{j=1}^{\infty}A_j\|_{S_p\left(\bigotimes \limits_{j=1}^{\infty}\Hi_j\right)}=\prod\limits_{j=1}^{\infty}\|A_j\|_{S_p(\Hi_j)}.\]
In particular, when $p=1$, for the trace class we additionally have 
\begin{equation}
\Tr\left(\bigotimes \limits_{j=1}^\infty A_j\right)=\prod\limits_{j=1}^{\infty}\Tr(A_j).\label{trinf2}\end{equation}	
\end{thm}%
\begin{proof}%
\noindent {\mbox{(i)}} Since the operators $A_j$ are compact and diagonal. They are also bounded and diagonal. Moreover, the convergence of  the net of partial products $\{\prod\limits_{j\in F}\|A_j\|_{\mathcal{L}(\Hi_j)}\}_{F\in\fou}$, implies that this net is also bounded, and by Theorem \ref{inftetr1} (i), the tensor $\bigotimes \limits_{j=1}^\infty A_j$ is bounded and diagonal. To see that this tensor is compact, it will be  enough to show that the net(sequence) of eigenvalues given by \eqref{eigeform8}  
 diverges to zero. But this follows from the first  inequality in \eqref{eignormiq1} and our assumption.\\

\noindent {\mbox{(ii)}} The fact that $\bigotimes \limits_{j=1}^{\infty}A_j$ is compact diagonal follows from (i), since $\|A_j\|_{\mathcal{L}(\Hi_j)}\leq \|A_j\|_{S_p(\Hi_j)}$.\\

On the ther hand, since the singular values of $\bigotimes \limits_{j=1}^{\infty}A_j$ can be obtained from \eqref{eigeform8}, for each finite set $F\in\fou$ we have a singular value of the form
\[\prod\limits_{j\in F}|\lambda_{k_{j}}^{(j)}|.\]


**Since every $A_j\in S_p(\Hi_j)$, is compact and diagonal we observe that the tensor is diagonal by (i). On the other hand, we enumerate the eigenvalues in the diagonalization \eqref{}, every finite set of subindices $\{k_1,k_2,\dots, k_n\}$ of $n$ elements it is denoted by $(\ell)_n$ we call  *** by ****we write
\[\lambda_{(\ell)}:=\lambda_{k_1}^{(1)}\cdots \lambda_{k_n}^{(n)}. \]
To see that $\bigotimes \limits_{j=1}^{\infty}A_j$ is compact, it will be enough to prove that $\lim_{}\lambda_{(\ell)_n}=0.$

We now observe that for a given element $\psi$ in the orthonormal basis  


We  note that the operator .. is well-defined and diagonal, indeed...\\
As in \eqref{} we...$F=\{j_1,\dots ,j_{\ell}\}$
\[\sum\limits_{F\in\fou}\lambda_{k_{j_1}}^{(j_1)}\cdots \lambda_{k_{j_{\ell}}}^{(j_{\ell})}\]


\noindent (iii) We now consider 
\end{proof}%
\begin{rem} In the context of infinite products, if the partial products of a sequence of positive terms go to zero, it is customary to say that the product diverges  to zero, that is why we keep that convention.\\% 

Regarding the condition for compacteness on Theorem \ref{inftetr1b} (i). We note that, if the sequence of partial products $\{\prod\limits_{j=1}^n\|A_j\|_{\mathcal{L}(\Hi_j)}\}_{n}$ is decreasing and diverges to zero. Then, analogously to Remark \ref{remprod3}, we can see that $\|A_j\|_{\mathcal{L}(\Hi_j)}\leq 1$ for every $j\ge 2$.  Hence, the net of partial products  $\{\prod\limits_{j\in F}\|A_j\|_{\mathcal{L}(\Hi_j)}\}_{F\in\fou}$ is also decreasing and diverge to zero.    
\end{rem}




\begin{defn} \label{diagdef1} Let $\Hi$ be a complex separable Hilbert space. Let $T$ be a bounded linear operator. 
 We say that $T$ is a {\em diagonal operator} if there exists an orthonormal basis $\{e_n:n\in\ene\}$ of $\Hi$ and a bounded sequence of complex numbers $\{a_n:n\in\ene\}$, such that  \[Tx=\sum\limits_{n=1}^{\infty}a_n\langle e_n, x\rangle e_n,\]
    for all $x\in \Hi .$%
\end{defn}%
We denote by $\mathcal{F}$ the family of finite subsets of $\ene$, and endow it with the order of inclusion. The set $\mathcal{F}$ is countable and becomes a directed set. We will consider some complex-valued nets defined on $\fou$ and they will be used  to enumerate the eigenvalues of an infinite tensor of operators. \\%

\begin{rem} \label{remprod3} We will point out some basic facts regarding infinite products that we will use in the next theorem. We observe the following: if  $G$ is a directed countable set and $\Gamma$ is the family of finite sets of $G$ ordered by inclusion, given a net $\{b_{\gamma}\}_{\gamma\in G}$ of positive real numbers, assume that the net of partial products  $\{\prod\limits_{\gamma\in F}b_{\gamma}\}_{F\in\Gamma}$ is increasing and bounded. Then, it converges to the supremum and $b_{\gamma}\geq 1$ for each $\gamma$ that is not a minimal point of $G$.\\ 

In fact, if $\gamma$ is not a minimal point of $G$, then, there exists $\alpha\in G$ such that $\alpha<\gamma$. Hence, $0<b_{\alpha}\leq b_{\alpha}b_{\gamma}$. Thus, $b_{\gamma}\geq 1$. \\

We also note that,  if  each term  $b_{\gamma}\geq 1$ for every no minimal point, then, the net of partial products $\{\prod\limits_{\gamma\in F}b_{\gamma}\}_{F\in\Gamma}$ is increasing, as this can be proved by induction. 

%The fact that $b_{\gamma}\geq 1$ follows from the inequality
%\[\prod\limits_{j=1}^{n+1}b_j=b_{n+1}\prod\limits_{j=1}^{n}b_j\geq %\prod\limits_{j=1}^{n}b_j.\]
%Hence $b_{n+1}\geq 1$ for every $n\geq 1$. On the other hand, if each term $b_n$ for $n\geq 2$, satisfies $b_n\geq 1$, then the sequence of partial products is increasing, as can be seen by induction.  \\

%In our theorem this says that a necessary condition for the norms of the operators is that $\|A_j\|_{\mathcal{L}(\Hi_j)}\geq 1$ for $j\geq 2$. 
\end{rem}%

**We can now establish our main theorem. We first obtain sufficient conditions for the compactness of infinite tensor products of compact diagonal operators.\\%

We consider  $\Hi_j$ to  be a complex separable Hilbert space for $j=1, 2, \dots \,\,$. We asume that  $A_j\in \mathcal{L}(\Hi_j)$  has a diagonal form as in Definition \ref{diagdef1}, for $j=1, 2, \dots \,$;  and that the net of partial products $\{\prod\limits_{j\in F}\|A_j\|_{\mathcal{L}(\Hi_j)}\}_{F\in\fou}$ is bounded. \\

For  each $k$, we write $A_k$ in a diagonal form 
\[A_kx=\sum\limits_{J=1}^{\infty}\lambda_j^{(k)}\langle e_j^{k},x\rangle e_j^{k}.\]
We denote by $\Lambda _k$ the orthonormal basis 
$\{e_j^{k}:j=1,2\dots\}$ of $\Hi_k$ and by $\Lambda$ the orthonormal basis of $\bigotimes \limits_{k=1}^{\infty}\Hi_k$  as in Theorem \ref{thmtesepar}. Let us fix the stabilizing sequence  $u_j$ by $u_n:=e_{1}^{(n)}$. We now observe that for a given  $\psi$ in  the orthonormal basis  $\Lambda$, there exists a finite set $\{k_1,\dots,k_n\}$ such that $\psi$ can be written in the form
\begin{equation}\label{psif1}
\psi=e_{k_1}^1\otimes e_{k_2}^2\otimes\cdots\otimes e_{k_n}^n\otimes u_{n+1}\otimes u_{n+2}\otimes\cdots  \, , 
\end{equation}%
noticing that some of the $e_{k_j}^j$'s can be equal to the respective $u_j$.\\%
We now define $\left(\bigotimes\limits_{k=1}^{\infty}A_k\right)(\psi)$ as an element of $\bigotimes \limits_{k=1}^{\infty}\Hi_k$, by
\begin{equation}\label{atensortilde}%
\left(\bigotimes\limits_{k=1}^{\infty}A_k\right)(\psi):=\left(\bigotimes\limits_{k=1}^{n}\widetilde{A}_k(\psi_k)\right) \otimes u_{n+1}\otimes u_{n+2}\otimes\cdots  \, ,
\end{equation}%
where 
\begin{equation} \label{psicases1}
    \widetilde{A}_k(\psi_k)= \begin{cases}A_k(\psi_k),\,\, & \mbox{ if }k=k_j\neq 1 \\
    \psi_k=e_1^k=u_k,\,\, & \mbox{ if }k=k_j=1.
    \end{cases}%
\end{equation}%
We note that the above tensor is well-defined. \\  
****

Thus, for each $\psi$ in the orthonormal basis $\Lambda$ we have defined 
$\left(\bigotimes\limits_{k=1}^{\infty}A_k\right)(\psi)$, according to formula \eqref{atensortilde} as an element of $\bigotimes \limits_{k=1}^{\infty}\Hi_k$. In the next theorem we  will establish the main properties of such mapping, starting with the boundedness for infinite tensors of diagonal operators as well as formulae for their norms. %
\begin{thm} \label{inftetr1} Let  $\Hi_j$  be a complex separable Hilbert space for $j=1, 2, \dots \,$. \\%

\noindent $(i)$ We assume that  $A_j\in \mathcal{L}(\Hi_j)$  has a diagonal form as in Definition \ref{diagdef1}, for $j=1, 2, \dots \,$;  and that the net of partial products $\{\prod\limits_{j\in F}\|A_j\|_{\mathcal{L}(\Hi_j)}\}_{F\in\fou}$ converges to zero. Then, the operator  $\bigotimes \limits_{j=1}^{\infty}A_j$ defined by $\left(\bigotimes\limits_{k=1}^{\infty}A_k\right)(\psi)$,  for each $\psi$ in the orthonormal basis $\Lambda$ as in \eqref{atensortilde}, is a compact diagonal operator.\\

%Moreover, if additionally the sequence of partial products  $\{\prod\limits_{j=1}^n\|A_j\|_{\mathcal{L}(\Hi_j)}\}_{F\in\fou}$ is increasing, then,  the operator norm of the tensor product $\bigotimes \limits_{j=1}^{\infty}A_j$ satisfies %
%\[\|\bigotimes %\limits_{j=1}^{\infty}A_j\|_{\mathcal{L}\left(\bigotimes \limits_{j=1}^{\infty}\Hi_j\right)}=\prod\limits_{j=1}^{\infty}\|A_j\|_{\mathcal{L}(\Hi_j)}. \]

\noindent $(ii)$ We asume that  $A_j\in \mathcal{L}(\Hi_j)$  has a diagonal form as in Definition \ref{diagdef1}, for $j=1, 2, \dots \,$;  and that the net of partial products $\{\prod\limits_{j\in F}\|A_j\|_{\mathcal{L}(\Hi_j)}\}_{F\in\fou}$ converges to zero. Then the adjoint of $\bigotimes \limits_{j=1}^{\infty}A_j$ is a compact diagonal operator with respect to the same orthonormal basis  and 
\[\left(\bigotimes \limits_{j=1}^{\infty}A_j\right)^*=\bigotimes \limits_{j=1}^{\infty}A_j^*.\]

\noindent {\mbox{(iii)}} Under the assumptions of (ii), we have that the operator 
\[\left(\bigotimes \limits_{j=1}^{\infty}A_j\right)^*\circ\left(\bigotimes \limits_{j=1}^{\infty}A_j\right),\]
is a bounded diagonal operator and 
\begin{equation}\label{r45} 
\left(\bigotimes \limits_{j=1}^{\infty}A_j\right)^*\circ\left(\bigotimes \limits_{j=1}^{\infty}A_j\right)=\bigotimes \limits_{j=1}^{\infty}(A_j^*A_j).\end{equation}
Consequently, the operator $\left|\bigotimes \limits_{j=1}^{\infty}A_j\right|$
 is also a diagonal operator and 
 \[\left|\bigotimes \limits_{j=1}^{\infty}A_j\right|=\bigotimes \limits_{j=1}^{\infty}|A_j|.\]%
\end{thm}

\begin{proof} (i)  We first observe that since the oprators $A_j$ are compact, they are also bounded diagonal and applying Theorem .. of \cite{DPR1}, we have that the infinite tensor $\bigotimes\limits_{j=1}^{\infty}A_j$ is also diagonal. To prove the compactness of this tensor, it will be enough to see that  its  eigenvalues are countable and converge to zero. Indeed, the eigenvalues can be obtained  from  \eqref{atensortilde}, denoting  by $F$ the set $\{j_1,j_2,\dots, j_{\ell}\}$ of indices $j=1,2,\dots, n$ such that $k_{j_i}\neq 1$, that is so that  $e_{k_{j_i}}^{j_i}\neq u_{k_{j_i}}$. We obtain 
\[ \left\langle \left(\bigotimes\limits_{j=1}^{\infty}A_j\right)(\psi), \psi \right\rangle = \]
\begin{align*} =&\left\langle\left(\bigotimes\limits_{j=1}^{n}\widetilde{A}_j(\psi_j)\right) \otimes u_{n+1}\otimes u_{n+2}\otimes\cdots  \, ,\left(\bigotimes\limits_{j=1}^{n}e_{k_j}^j\right) \otimes u_{n+1}\otimes u_{n+2}\otimes\cdots \right \rangle\\\nonumber
=&\prod\limits_{i=1}^{\ell}\lambda_{k_{j_i}}^{(j_i)}\,\,  \\\nonumber
=&\prod\limits_{j\in F}\lambda_{k_{j}}^{(j)}\,  . 
\end{align*}%
Hence%
\begin{equation}\label{psidiag}
    \left\langle\left(\bigotimes\limits_{j=1}^{\infty}A_j\right)(\psi), \psi \right\rangle =\prod\limits_{j\in F}\lambda_{k_{j}}^{(j)}\, .
\end{equation}%
The net $\prod\limits_{j\in F}\lambda_{k_{j}}^{(j)}$ constitute the family of eigenvalues of the infinite tensor
 $\bigotimes\limits_{j=1}^{\infty}A_j$. We now observe that since the  net of partial products $\{\prod\limits_{j\in F}\|A_j\|_{\mathcal{L}(\Hi_j)}\}_{F\in\fou}$ converges to zero and%
\begin{equation}\label{eignormiq1}
  \prod\limits_{j\in F}|\lambda_{k_{j}}^{(j)}|\leq \prod\limits_{j\in F}\|A_j\|_{\mathcal{L}(\Hi_j)}. \end{equation}%
Then, the net of eigenvalues  $\{\prod\limits_{j\in F}\lambda_{k_{j}}^{(j)}\}_{F\in\fou}$ also converges to zero. On the other hand, it is clear that the direct set $\mathcal{F}$ is countable. Therefore, the tensor  $\bigotimes\limits_{k=1}^{\infty}A_k$ is compact.\\

(ii) We just need to prove that the adjoint of $\bigotimes\limits_{k=1}^{\infty}A_k$
******************

  
Therefore, we have also obtained that the operator $\bigotimes\limits_{k=1}^{\infty}A_k$ is a bounded diagonal operator with eigenvalues
\begin{equation}\label{eigeform8}%
    \prod\limits_{j\in F}\lambda_{k_{j}}^{(j)},
\end{equation}%
for each $F\in\fou$. This proves the first part of (i).\\%


Regarding the formula for the operator norm, we first note that if the sequence of partial products  $\{\prod\limits_{j=1}^n\|A_j\|_{\mathcal{L}(\Hi_j)}\}_{n}$ is increasing, so  is 
$\{\prod\limits_{j\in F}\|A_j\|_{\mathcal{L}(\Hi_j)}\}_{F\in\fou}$. Indeed, by Remark \ref{remprod3}, we have that $\|A_j\|_{\mathcal{L}(\Hi_j)}\geq 1$ for $j\geq2$. Hence, the net of partial products $\{\prod\limits_{j\in F}\|A_j\|_{\mathcal{L}(\Hi_j)}\}_{F\in\fou}$ is increasing and bounded. \\

For $F\in\fou$ we note that if $n=\max F$, then $F\subset \widetilde{F}:=\{1,...,n\}$, and since  $\bigotimes\limits_{k=1}^{\infty}A_k$ is diagonal, using \eqref{psidiag} we have that
\begin{align*}
\|\bigotimes \limits_{j=1}^{\infty}A_j\|_{{\mathcal{L}}\left(\bigotimes \limits_{j=1}^{\infty}\Hi_j\right)}&=\sup\limits_{\psi\in\Lambda}\left |\left\langle \left(\bigotimes\limits_{k=1}^{\infty}A_k\right)(\psi),\psi\right\rangle\right|\\
&=\sup\limits_{F\in\fou}\prod\limits_{j\in F}|\lambda_{k_{j}}^{(j)}|\\
&\leq \sup\limits_{\widetilde{F}\in\fou}\prod\limits_{j\in \widetilde{F}}\|A_j\|_{{\mathcal{L}}(\Hi_j)}\\
&=\sup\limits_{n}\prod\limits_{j=1}^{n}\|A_j\|_{{\mathcal{L}}(\Hi_j)}\\
&=\prod\limits_{j=1}^{\infty}\|A_j\|_{{\mathcal{L}}(\Hi_j)},
\end{align*}
where we have considered the elements $\psi$ of the orthonormal basis $\Lambda$ written in the form \eqref{psif1} to evaluate the inner product. The last identity comes from the fact that the partial products form an increasing 
 net. This shows that 
 \[\|\bigotimes \limits_{j=1}^{\infty}A_j\|_{{\mathcal{L}}\left(\bigotimes \limits_{j=1}^{\infty}\Hi_j\right)}\leq \prod\limits_{j=1}^{\infty}\|A_j\|_{{\mathcal{L}}(\Hi_j)}.\]
 We now prove the opposite inequality. We will use the fact that  $\|A_j\|_{{\mathcal{L}}(\Hi_j)}=\sup\limits_k|\lambda_{k}^{(j)}|, $ for all $j$. Hence, for every $n$ we have%
 \begin{align*}
     \prod\limits_{j=1}^{n}\|A_j\|_{{\mathcal{L}}(\Hi_j)}&= \prod\limits_{j=1}^{n}\sup\limits_k|\lambda_{k}^{(j)}|\\
     &= \sup\limits_{k}\prod\limits_{j=1}^{n} |\lambda_{k}^{(j)}|.
      \end{align*}
We now choose for each $k$ a vector $\psi_k\in\Lambda$ defined by
\[\psi_k:=e_{k}^1\otimes e_{k}^2\otimes\cdots\otimes e_{k}^n\otimes u_{n+1}\otimes u_{n+2}\otimes\cdots  \, .\]
Hence 
\begin{align*}
     \sup\limits_{k}\prod\limits_{j=1}^{n} |\lambda_{k}^{(j)}|&= \sup\limits_{k} \left |\left\langle \left(\bigotimes\limits_{i=1}^{\infty}A_i\right)(\psi_k),\psi_k\right\rangle\right|\\
     &\leq \|\bigotimes \limits_{j=1}^{\infty}A_j\|_{{\mathcal{L}}\left(\bigotimes \limits_{j=1}^{\infty}\Hi_j\right)}.
      \end{align*}%
This concludes the proof of (i).\\%

\noindent (ii) Since each $A_j$ is a bounded diagonal operator, then  every  $A_j^*$ is a bounded diagonal operator. Moreover, since $\|A_j^*\|_{\mathcal{L}(\Hi_j)}=\|A_j\|_{\mathcal{L}(\Hi_j)}$ we get that the net of partial products $\{\prod\limits_{j\in F}\|A_j^*\|_{\mathcal{L}(\Hi_j)}\}_{F\in\fou}$ is bounded.  Hence, 
 the tensor product of operators $\bigotimes \limits_{j=1}^{\infty}A_j^*$ is well-defined and bounded diagonal  by (i). On the other hand, the tensor product $\bigotimes \limits_{j=1}^{\infty}A_j$ being also bounded and diagonal  by (i), its adjoint is defined as a bounded operator. Moreover, it is not difficult to see that the adjoint of a bounded diagonal operator with eigenvalues $\{\lambda_k\}_k$, is a bounded diagonal operator with respect to the same orthonormal basis and eigenvalues $\{\overline{\lambda_k}\}_k$. Therefore, for every element $\psi$ of the orthonormal basis $\Lambda$, by using the formula for the eigenvalues \eqref{eigeform8}, we have
 \begin{align*} \left(\bigotimes \limits_{j=1}^{\infty}A_j\right)^*(\psi)=&\prod\limits_{j\in F}\overline{\lambda_{k_{j}}^{(j)}}\,\,\psi\\
 =&\prod\limits_{j\in F}A_j^*(\psi_j)\\
 =&\left(\bigotimes \limits_{j=1}^{\infty}A_j^*\right)(\psi).
  \end{align*}
Thus
\[\left(\bigotimes \limits_{j=1}^{\infty}A_j\right)^*=\bigotimes \limits_{j=1}^{\infty}A_j^*,\]
completing the proof of (ii).\\

\noindent (iii) For the finite case $\bigotimes \limits_{j=1}^{N}A_j$, even just with the boundedness assumption on each operator, it is not difficult to see that  $\left(\bigotimes \limits_{j=1}^{N}A_j\right)^*=\bigotimes \limits_{j=1}^{N}A_j^*$. For the infinite countable case, we use the diagonality assumption. Then $\bigotimes \limits_{j=1}^{\infty}A_j^*$ is a bounded diagonal operator with respect to the same orthornomal basis  as  $\bigotimes \limits_{j=1}^{\infty}A_j$. Let $\psi\in\Lambda$ be as in \eqref{psif1}.  Then, by \eqref{atensortilde} applied to $\bigotimes \limits_{j=1}^{\infty}A_j$ and $\bigotimes \limits_{j=1}^{\infty}A_j^*$ (using (ii)), we have  
\begin{align*}
\left(\bigotimes \limits_{j=1}^{\infty}A_j^*\right)\circ\left(\bigotimes \limits_{j=1}^{\infty}A_j\right)(\psi)&=\left(\bigotimes \limits_{j=1}^{\infty}A_j^*\right)\left(\left(\bigotimes \limits_{j=1}^{\infty}A_j\right)(\psi)\right)\\
&=\left(\bigotimes\limits_{j=1}^{\infty}A_j^*\right)\left(\left(\bigotimes\limits_{k=1}^{n}\widetilde{A}_k(\psi_k)\right) \otimes u_{n+1}\otimes u_{n+2}\otimes\cdots\right)\\
&=\left(\bigotimes\limits_{j=1}^{\infty}A_j^*\right)\left(\left(\bigotimes\limits_{j=1}^{n}(\lambda_{k_j}^{(j)}e_{k_j}^j)\right) \otimes u_{n+1}\otimes u_{n+2}\otimes\cdots\right)\\
&=\left(\bigotimes\limits_{j=1}^{n}|\lambda_{k_j}^{(j)}|^2e_{k_j}^j\right) \otimes u_{n+1}\otimes u_{n+2}\otimes\cdots\\
&=\left(\bigotimes\limits_{j=1}^{n}(A_j^*A_j)(e_{k_j}^j)\right) \otimes u_{n+1}\otimes u_{n+2}\otimes\cdots\\
&=\left(\bigotimes \limits_{j=1}^{\infty}(A_j^*A_j)\right)(\psi).
\end{align*} 
Thus, $\left(\bigotimes \limits_{j=1}^{\infty}A_j^*\right)\circ\left(\bigotimes \limits_{j=1}^{\infty}A_j\right)(\psi)$ is a bounded diagonal operator and \eqref{r45} holds. It is clear that this operator is also positve definite and 
\[\left|\bigotimes \limits_{j=1}^{\infty}A_j\right|=\bigotimes \limits_{j=1}^{\infty}|A_j|,\]%
completing the proof.
\end{proof}



\noindent $(i)$ Let $1\leq p<\infty$, we assume that  $A_j\in S_p(\Hi_j)$  has a diagonal form as in Definition \ref{diagdef1}, for $j=1, 2, \dots \,$;  and that the net of partial products $\{\prod\limits_{j\in F}\|A_j\|_{\mathcal{L}(\Hi_j)}\}_{F\in\fou}$ converges.


*******
We now explain how to define infinite tensor products of operators in Schatten-von Neumann classes. \\

Let  $\Hi_j$  be a complex separable Hilbert space for $j=1, 2, \dots \,$.  Let $1\leq p<\infty$, we consider operators  $A_j\in S_p(\Hi_j)$, for $j=1,2,\dots$. First, we will associate to each finite tensor  $\bigotimes \limits_{j=1}^{N}A_j$ defined on $\bigotimes \limits_{j=1}^{N}\Hi_j$, a tensor defined on the fixed Hilbert space $\bigotimes \limits_{j=1}^{\infty}\Hi_j$. Indeed, given $N\in\ene$, we identify the tensor $\bigotimes \limits_{j=1}^{N}A_j$ that belongs to $S_p\left(\bigotimes \limits_{j=1}^{N}\Hi_j\right)$, with the tensor 
\begin{equation}\label{opan}
    A_1\otimes A_2\otimes\cdots\otimes A_N\otimes I_{N+1}\otimes I_{N+2}\otimes\cdots=A_1\otimes A_2\otimes\cdots\otimes A_N\otimes\bigotimes \limits_{j=N+1}^{\infty}I_{j},
\end{equation}
where $I_j$ is the identity operator on $\Hi_j$.\\

We note that \eqref{opan} defines indeed an operator on $\bigotimes \limits_{j=1}^{\infty}\Hi_j$. \\

We denote
\[B:=\bigotimes \limits_{j=1}^{N}A_j\otimes \bigotimes \limits_{j=N+1}^{\infty}I_j,\]
and $\mathcal{M}_N:=\bigotimes \limits_{j=1}^{N}\Hi_j$. Since  $\mathcal{M}_N$ is a closed subspace of $\bigotimes \limits_{j=1}^{\infty}\Hi_j$, we can define its orthogonal projection $P_{\mathcal{M}_N}:\bigotimes \limits_{j=1}^{\infty}\Hi_j\rightarrow \mathcal{M}_N$. We observe that 
\[B=\bigotimes \limits_{j=1}^{N}A_j\circ P_{\mathcal{M}_N}. \]
Then
\[\|B\|_{S_p\left(\bigotimes\limits_{j=1}^{\infty}\Hi_j\right)}=\prod\limits_{j=1}^{N}\|A_j\|_{S_p(\Hi_j)}.\]



$A\left|_{\widetilde{\bigotimes\limits_{j=1}^{N}}\Hi_j}\right.$
IT IS clear that ...are isometrically isomorphic.\\


Regarding the compactness of an infinite tensor we obtain the following result that will also help to prepare our main result for the characterization of infinite tensors in Schatten-von Neumann classes.

\begin{thm} \label{main1aax} Let  $\Hi_j$  be a complex separable Hilbert space for $j=1, 2, \dots \,$. If  $A_j\in \mathcal{L}(H_j)$, 
 compact for all $j=1,2,\dots$. Then, each tensor $\widetilde{\bigotimes \limits_{j=1}^{N}}A_j$ is a compact operator on $\bigotimes\limits_{j=1}^{\infty}\Hi_j$.
 The  sequence of tensor products $\{\widetilde{\bigotimes \limits_{j=1}^{N}}A_j\}_{N\in\ene}$ is convergent in $\mathcal{L}\left(\bigotimes\limits_{j=1}^{\infty}\Hi_j\right)$ if and only if $\prod\limits_{j=1}^{\infty}\|A_j\|_{\LL(\Hi_j)}$ converges.\\

 If that is the case $\bigotimes \limits_{j=1}^{\infty}A_j$ is a compact operator on $\bigotimes\limits_{j=1}^{\infty}\Hi_j$ and
 \[\|\bigotimes \limits_{j=1}^{\infty}A_j\|=\prod\limits_{j=1}^{\infty}\|A_j\|_{\LL(\Hi_j)}.\]
\end{thm}
\begin{proof} We first prove that $\widetilde{\bigotimes \limits_{j=1}^{N}}A_j$ is compact. We denote by $B$ the open unit ball of the Hilbert space $\bigotimes \limits_{j=1}^{\infty}H_j$. We need to prove that $\left(\widetilde{\bigotimes \limits_{j=1}^{N}}A_j\right)(B)$ is relatively compact in $\bigotimes \limits_{j=1}^{\infty}H_j$. Indeed,  since the range $Ran (\widetilde{\bigotimes \limits_{j=1}^{N}}A_j)$ is contained in $\widetilde{\bigotimes \limits_{j=1}^{N}}H_j$, and $\widetilde{\bigotimes \limits_{j=1}^{N}}H_j$ is isometrically isomorphic to $\bigotimes \limits_{j=1}^{N}H_j$. Then, the conclusion  follows from  the fact that $\left(\widetilde{\bigotimes \limits_{j=1}^{N}}A_j\right)(B)$ is homeomorphic to $\left(\bigotimes \limits_{j=1}^{N}A_j\right)(B)$ and the fact that $\bigotimes \limits_{j=1}^{N}A_j$ is a compact operator. \\

As for the second part of the statement, we observe that for $M,N\in\ene$ with $M>N$, we have%
\begin{equation}\label{prte4}
\|\bigotimes \limits_{j=1}^{M}A_j-\bigotimes \limits_{j=1}^{N}A_j\|_{\LL\left(\bigotimes\limits_{j=1}^{N}\Hi_j\right)}=\|\bigotimes \limits_{j=N+1}^{M}A_j\|_{\LL\left(\bigotimes\limits_{j=1}^{N}\Hi_j\right)}=\prod\limits_{j=N+1}^{M}\|A_j\|_{\LL(\Hi_j)}.%
\end{equation}%
Therefore, $\{\widetilde{\bigotimes \limits_{j=1}^{N}}A_j\}_{N\in\ene}$ is a Cauchy sequence in $\LL\left(\bigotimes\limits_{j=1}^{\infty}\Hi_j\right)$ if and only if $\prod\limits_{j=1}^{\infty}\|A_j\|_{\LL(\Hi_j)}$ is a Cauchy sequence in $\er$. Hence, the sequence of tensor products $\{\widetilde{\bigotimes \limits_{j=1}^{N}}A_j\}_{N\in\ene}$ is convergent in $\mathcal{L}\left(\bigotimes\limits_{j=1}^{\infty}\Hi_j\right)$ if and only if $\prod\limits_{j=1}^{\infty}\|A_j\|_{\LL(\Hi_j)}$ converges.\\

Now if $\lim_N\widetilde{\bigotimes \limits_{j=1}^{N}}A_j=\bigotimes \limits_{j=1}^{\infty}A_j$ in $\LL\left(\bigotimes\limits_{j=1}^{\infty}\Hi_j\right)$. Then,  $\bigotimes\limits_{j=1}^{\infty}A_j$ is compact since it is the limit of compact operators with respect to the operator norm, and the ideal of compact operators is closed. By the continuity of the norm  in $\LL\left(\bigotimes\limits_{j=1}^{\infty}\Hi_j\right)$, we get the last conclusion of the theorem.
\end{proof}
We can now state below some immediate consequences for the adjoint of infinite tensors of compact operators.
\begin{cor}\label{adjo3} Let  $\Hi_j$  be a complex separable Hilbert space for $j=1, 2, \dots \,$. If  $A_j\in \mathcal{L}(H_j)$, 
 is compact for all $j=1,2,\dots$ and  $\prod\limits_{j=1}^{\infty}\|A_j\|_{\LL(\Hi_j)}$ converges. Then
 
 \noindent {\mbox{(i)}} The adjoint of $\bigotimes \limits_{j=1}^{\infty}A_j$ is a compact
operator and 
\[\left(\bigotimes \limits_{j=1}^{\infty}A_j\right)^*=\bigotimes \limits_{j=1}^{\infty}A_j^*.\]

 \noindent {\mbox{(ii)}} The operator 
\[\left(\bigotimes \limits_{j=1}^{\infty}A_j\right)^*\circ\left(\bigotimes \limits_{j=1}^{\infty}A_j\right),\]
is  compact   and 
\begin{equation}\label{r45s} 
\left(\bigotimes \limits_{j=1}^{\infty}A_j\right)^*\circ\left(\bigotimes \limits_{j=1}^{\infty}A_j\right)=\bigotimes \limits_{j=1}^{\infty}(A_j^*A_j).
\end{equation}
Consequently, the operator $\left|\bigotimes \limits_{j=1}^{\infty}A_j\right|$
 is compact and 
 \[\left|\bigotimes \limits_{j=1}^{\infty}A_j\right|=\bigotimes \limits_{j=1}^{\infty}|A_j|.\]%
\end{cor}
\begin{proof} (i) Since each $A_j$ is a compact operator in  $\mathcal{L}(\Hi_j)$, then  every  $A_j^*$ is compact. Moreover, since $\|A_j^*\|_{\mathcal{L}(\Hi_j)}=\|A_j\|_{\mathcal{L}(\Hi_j)}$ we get that $\prod\limits_{j=1}^{\infty}\|A_j^*\|_{\LL(\Hi_j)}$ converges. Hence $\bigotimes \limits_{j=1}^{\infty}A_j^*$ is compact. Now the conclusion follows from the continuity of the convolution and the convergence \[\left(\bigotimes \limits_{j=1}^{N}A_j^\right)^*=\bigotimes \limits_{j=1}^{N}A_j^*\longrightarrow \left(\bigotimes \limits_{j=1}^{\infty}A_j^\right)^*\]
in $\LL\left(\bigotimes \limits_{j=1}^{\infty}A_j\right)$.\\

\noindent (ii) We first observe that  

\[\left(\bigotimes \limits_{j=1}^{\infty}A_j\right)^*\circ\left(\bigotimes \limits_{j=1}^{\infty}A_j\right)=\left(\bigotimes \limits_{j=1}^{\infty}A_j^*\right)\circ\left(\bigotimes \limits_{j=1}^{\infty}A_j\right).\]
On the other hand, for all $N$ we have
\[\left(\bigotimes \limits_{j=1}^{N}A_j^*\right)\circ\left(\bigotimes \limits_{j=1}^{N}A_j\right)=\bigotimes \limits_{j=1}^{N}(A_j^*A_j).\]
Since $A_j^*A_j\in \LL(\Hi_j)$ is compact for all $J$, then we get a sequence compact tensors $\{\bigotimes \limits_{j=1}^{N}A_j^*A_j\}_{N\in\ene}$ that converges to a compact operator $\bigotimes \limits_{j=1}^{\infty}A_j^*A_j$.
 Indeed, looking at the norms of the terms in the sequence 
\begin{align*} \left\|\bigotimes \limits_{j=1}^{N}A_j^*A_j\right\|_{\LL\left(\bigotimes\limits_{j=1}^{\infty}\Hi_j\right)}=&\prod\limits_{j=1}^{N}\|A_j^*A_j\|_{\LL(\Hi_j)}\\
=&\prod\limits_{j=1}^{N}\|A_j\|_{\LL(\Hi_j)}^2\\
 =&\left(\prod\limits_{j=1}^{N}\|A_j\|_{\LL(\Hi_j)}\right)^2.
\end{align*}
Since 
$\prod\limits_{j=1}^{\infty}\|A_j\|_{\LL(\Hi_j)}$ converges, the same holds for $\left(\prod\limits_{j=1}^{N}\|A_j\|_{\LL(\Hi_j)}\right)^2$ as $N\rightarrow\infty.$ Hence, $\bigotimes \limits_{j=1}^{\infty}A_j^*A_j$ is compact with norm 
\[\left\|\bigotimes \limits_{j=1}^{\infty}A_j^*A_j\right\|_{\LL\left(\bigotimes\limits_{j=1}^{\infty}\Hi_j\right)}=\left(\prod\limits_{j=1}^{\infty}\|A_j\|_{\LL(\Hi_j)}\right)^2.\]
To conclude the proof we only need to show that 
\[\lim\limits_{N\rightarrow \infty}\left(\bigotimes \limits_{j=1}^{N}A_j^*\right)\circ\left(\bigotimes \limits_{j=1}^{N}A_j\right)=\left(\bigotimes \limits_{j=1}^{\infty}A_j^*\right)\circ\left(\bigotimes \limits_{j=1}^{\infty}A_j\right),\]
with respect to the norm in $\LL\left(\bigotimes\limits_{j=1}^{\infty}\Hi_j\right)$. But this follows from  the fact that $\LL\left(\bigotimes\limits_{j=1}^{\infty}\Hi_j\right)$ is a Banach algebra and 
\[\lim\limits_{N\rightarrow \infty}\left(\bigotimes \limits_{j=1}^{N}A_j^*\right)=\bigotimes \limits_{j=1}^{\infty}A_j^*\]
and
\[\lim\limits_{N\rightarrow \infty}\left(\bigotimes \limits_{j=1}^{N}A_j\right)=\bigotimes \limits_{j=1}^{\infty}A_j.\]
By the continuity of the operation $(A,B)\rightarrow AB=A\circ B$, on the Banach algebra $\LL\left(\bigotimes\limits_{j=1}^{\infty}\Hi_j\right)$ we can conclude the proof  of \eqref{r45s}.\\

The remainder part of (ii) follows from***

\end{proof}




Gauge field theory coherent states (GCS): IV. Infinite tensor product and thermodynamical limit

Hydrogenic atom:
any atom or ion with a single valence electron 
